\newif\ifcomments     %% include author discussion
\newif\ifanonymous    %% include author identities
\newif\ifextended     %% include appendix
\newif\ifsubmission   %% prepare the submitted version
\newif\ifpublic       %% version available for posting / final version

\commentsfalse        %% toggle comments here
\extendedfalse

\submissionfalse      %% at most one of these must be true (neither for draft version)
\publictrue          %% but if you want to see comments, these should both be off


% If we are going to make a version public, i.e. on arXiv, we should 
% make sure that there are no comments and our names are on it.
\ifpublic
\submissionfalse
\commentsfalse
\anonymousfalse
\fi

%% If we are submission, make sure there are no comments and our names 
%% are NOT on it.
\ifsubmission
\publicfalse
\commentsfalse
\anonymoustrue
\fi


\PassOptionsToPackage{dvipsnames,svgnames,x11names}{school}
% \documentclass[acmsmall,screen=true,
% \ifpublic review=false\else,review=true\fi
% ,anonymous=\ifanonymous true\else false\fi]{acmart}
\documentclass[acmsmall]{acmart}
\citestyle{acmauthoryear}
%%% The following is specific to POPL '24 and the paper
%%% 'Internalizing Indistinguishability with Dependent Types'
%%% by Yiyun Liu, Jonathan Chan, Jessica Shi, and Stephanie Weirich.
%%%
\setcopyright{rightsretained}
\acmDOI{10.1145/3632886}
\acmYear{2024}
\copyrightyear{2024}
\acmSubmissionID{popl24main-p229-p}
\acmJournal{PACMPL}
\acmVolume{8}
\acmNumber{POPL}
\acmArticle{44}
\acmMonth{1}
\received{2023-07-11}
\received[accepted]{2023-11-07}

\usepackage{natbib}
\citestyle{acmauthoryear}
\usepackage{ottalt}
\usepackage{mathpartir}
\usepackage{supertabular}
\usepackage[para]{footmisc}
\usepackage{amsmath}

\usepackage{draft}  % local

\ifcomments
\newnote{scw}{blue} % Stephanie Weirich
\newnote{yl}{purple} % Yiyun Liu
\newnote{jc}{olive} % Jonathan Chan
\newnote{jws}{orange} % Jessica Shi
\else
\newcommand{\scw}[1]{}
\newcommand{\yl}[1]{}
\newcommand{\jc}[1]{}
\newcommand{\jws}[1]{}
\fi


\usepackage{color}
\inputott{Qualitative-rules}

\usepackage{url}
\usepackage{listings}
\lstset{ % this is from Harry's paper, feel free to replace w other settings
  basicstyle=\sffamily,
  columns=flexible,
  mathescape=true,
  breaklines=false,
  morekeywords={data,do,where,return,type,case,of,if,then,else,let,in,levels,constraints,refl,transp},
  %stringstyle=\color{string},
  %morestring=*[d]{"}, % chktex 18
  %morestring=*[d]{'},
  literate={*>}{{\kl[applicative]{$*\rangle$}}}2
  {<*>}{{\kl[applicative]{$\langle * \rangle$}}}3
  {->}{{$\mathrel{\rightarrow}$}}2
  {\\}{{$\lambda$}}1
  {==>}{{$\mathrel{\implies}$}}3
  {=>}{{$\mathrel{\Rightarrow}$}}2
  {<|>}{{$\langle | \rangle$}}2
  {|-}{{$\vdash$}}1
  {pure}{\kl[applicative]{\textsf{pure}}}4
  {fmap}{{$\fmap$}}3
  {fst}{{$\pi_1\;$}}1
  {snd}{{$\pi_2\;$}}1
  {~}{{$\mathbin{\sim}$}}1
  {'}{{$^\prime$}}0
}
\def\lstinlinecol{\lstinline[columns=fullflexible]}

\usepackage[noend]{algpseudocode}
\usepackage{xspace}

\newcommand\suppress[1]{}
\newcommand{\name}{DCOI\xspace}
\newcommand{\Name}{DCOI\xspace}
\newcommand{\longname}{Dependent Calculus of Indistinguishability\xspace}
\newcommand{\subst}[3]{#1\{#2/#3\}}
\newcommand{\dotv}[3]{\href{#1}{\texttt{#2}}{\texttt{:#3}}}
\newcommand{\dotpi}[2]{\href{#1/#2}{\texttt{#2}}}

% use this before footnote to stack footnote mark above punctuation
\newlength{\punctwidth}
\newcommand{\punctstack}[1]{#1%
  \settowidth{\punctwidth}{#1}%
  \hspace*{-\the\punctwidth}%
}

%%
%% end of the preamble, start of the body of the document source.
\begin{document}
\title{Internalizing Indistinguishability with Dependent Types}

\author{Yiyun Liu}
\orcid{0009-0006-8717-2498}
\affiliation{%
  \institution{University of Pennsylvania}
  \country{USA}
}
\email{liuyiyun@seas.upenn.edu}

\author{Jonathan Chan}
\orcid{0000-0003-0830-3180}
\affiliation{%
  \institution{University of Pennsylvania}
  \country{USA}
}
\email{jcxz@seas.upenn.edu}

\author{Jessica Shi}
\orcid{0000-0002-1507-1122}
\affiliation{%
  \institution{University of Pennsylvania}
  \country{USA}
}
\email{jwshi@seas.upenn.edu}

\author{Stephanie Weirich}
\orcid{0000-0002-6756-9168}
\affiliation{%
  \institution{University of Pennsylvania}
  \country{USA}
}
\email{sweirich@seas.upenn.edu}


\begin{CCSXML}
<ccs2012>
   <concept>
       <concept_id>10003752.10003790.10011740</concept_id>
       <concept_desc>Theory of computation~Type theory</concept_desc>
       <concept_significance>500</concept_significance>
       </concept>
 </ccs2012>
\end{CCSXML}

\ccsdesc[500]{Theory of computation~Type theory}

%%
%% By default, the full list of authors will be used in the page
%% headers. Often, this list is too long, and will overlap
%% other information printed in the page headers. This command allows
%% the author to define a more concise list
%% of authors' names for this purpose.
% \renewcommand{\shortauthors}{Liu et al.}

%%
%% The abstract is a short summary of the work to be presented in the
%% article.
\begin{abstract}
  In type systems with dependency tracking, programmers can assign an ordered
  set of levels to computations and prevent information flow from high-level
  computations to the low-level ones. The key notion in such systems is
  \emph{indistinguishability}: a definition of program equivalence that takes
  into account the parts of the program that an observer may depend on.  In
  this paper, we investigate the use of dependency tracking in the context of
  dependently-typed languages. We present the Dependent Calculus of
  Indistinguishability (DCOI), a system that adopts indistinguishability as
  the definition of equality used by the type checker. DCOI also internalizes
  that relation as an observer-indexed propositional equality type, so that
  programmers may reason about indistinguishability within the language. Our
  design generalizes and extends prior systems that combine dependency
  tracking with dependent types and is the first to support conversion and
  propositional equality at arbitrary observer levels.  We have proven type
  soundness and noninterference theorems for DCOI and have developed a
  prototype implementation of its type checker.
\end{abstract}

% Do we get to pick arbitrary keywords, or are they from a fixed ACM official list?
% We can pick arbitrary keywords, but don't need to do so until the final version.
% We will eventually have to pick from ACM's subject classification list, but not until later.
\keywords{Modes, Dependent Types, Coq, Formalization}

\maketitle

\section{Introduction}\label{sec:intro}
% In a type system that supports dependency tracking, the notion of
% indistinguishability can be used to equate programs that are
% syntactically distinct but appear identical up to erasure at a
% specific observer level.
Dependency tracking is a static analysis that determines how computations
depend on their inputs. Type systems that support dependency tracking assign
levels drawn from some partial order to computations, and prevent the flow of
information from higher level computations to lower ones.

For example, consider the type signature of a function \lstinline{f}
as follows, where the levels $[[Low]]$ (for low) and $[[High]]$ (for high) are ordered
$[[Low]] < [[High]]$. 
\begin{lstlisting}
    f :$^{[[Low]]}$ Nat$^{[[High]]}$ ->  Nat
\end{lstlisting}
% Consider a function variable with the
% following type signature in a type system that supports dependency tracking.
The type signature reads that \lstinline{f} is a function
that takes a natural number typed at level $[[High]]$ and returns a natural that is
typed at level $[[Low]]$. Because $[[Low]] < [[High]]$,
the function $[[f]]$ may not use its argument and must behave like a constant
function.

Dependency tracking is a general feature of type systems and has a variety of
applications~\citep{abadi1999core}, most frequently information flow
analysis~\citep{denning1977infoflow, sabelfeld2003language}, but also partial
evaluation~\citep{hatcliff_danvy_1997} and staged
programming~\citep{sheard1995staged, taha2000metaml}. In the context of
dependent type systems, it is related to relevance
tracking~\citep{debruijn1994automath,miquel:icc,mishra2008erasure,choudhury:ddc},
two-level type theory~\citep{annenkov2023two}, and the phase distinction in the
ML module system~\citep{sterling-harper:phase-distinction}.

Levels used for dependency tracking take on different meanings depending on
the specific application. For example, in relevance tracking, the dependency
levels $[[Low]]$ and $[[High]]$ can represent relevant and irrelevant
computations, respectively, where irrelevant computations can be erased prior to
run time. For information flow, the dependency levels correspond to clearance
levels and only computations typed at level $[[High]]$ can access classified
information.
% The dependency levels $[[Low]]$ and
% $[[High]]$ can be interpreted as security levels or relevance levels.
% Depending on the specific application, the set of levels can be drawn
% from a richer structure such as a lattice.
For each application, dependency levels are drawn from a predetermined
partially ordered set or a richer structure such as a lattice.

% In both cases, the type system
% prevents the flow of information from the higher level to the lower
% level.
When reasoning about program equivalence in a system that supports dependency
tracking, the key idea is \emph{indistinguishability}: an equivalence relation
that is indexed by an observer level~\citep{sabelfeld2003language}. This
relation defines program equivalence from the point of view of an
observer---all parts of the program that are unobservable at that level can be
safely equated because the system prevents the observer from depending on
those values.

Dependently-typed languages support intrinsic reasoning about program
equivalence. Calculi such as
PTS~\citep{barendregt:lambda-calculi-with-types}, CIC~\citep{CIC}, and
MLTT~\citep{Martin-Lof-1973} characterize what it means for programs to be
equal and use this definition of equality for type conversion.  In these
systems, two terms are typically definitionally equal when they are
$\beta$-equivalent or $\beta\eta$-equivalent. Furthermore, the propositional
equality type internalizes the equality judgment as a mechanism for intrinsic
reasoning about program equivalence. In CIC and MLTT, the type $[[a ~ b : A]]$
asserts that the terms $[[a]]$ and $[[b]]$ (of type $[[A]]$) are equal.  If
this proposition is provable, we say that $[[a]]$ and $[[b]]$ are
propositionally equal. In dependent type systems, more terms are
propositionally equal than definitionally equal: for example, one can use
induction on naturals to prove $[[x + 0 ~ x : Nat]]$ even when $[[x+0]]$
and $[[x]]$ are not definitionally equal.  By internalizing definitional
equality as the equality type, programmers gain a powerful tool for reasoning
about program equivalence.

However, the power of this tool depends on the definition of equivalence in
the first place. In the context of dependency tracking, it is sound to equate
programs that are not necessarily $\beta\eta$-equivalent as long as their
differences are not observable. To explore this idea, we present the \longname (\name{})\footnote{pronounced ``decoy''}, 
a dependently-typed calculus that uses
\emph{indistinguishability} as its definition of equality.

\Name{} can type check programs that would be rejected by CIC or MLTT. For
example, given a function \lstinline{P} with the signature
\lstinline{P :$^{[[Low]]}$ Bool$^{[[High]]}$ ->  Type},
% $ useless comment to fix emacs's broken syntax highlighting
we should be able to safely convert between the types \lstinline{P True} and
\lstinline{P False}. During conversion, the type checker can skip the comparison of
\lstinline{True} and \lstinline{False} because the signature of \lstinline{P}
tells us that \lstinline{P} cannot use its argument in a way that is
observable at level $[[Low]]$.

\medskip
In exploring the design of DCOI, this paper makes the following contributions.
\begin{itemize}
\item When using indistinguishability as definitional equality, there is the
  question of what observer level should be used. While it is tempting to use
  a fixed observer level for type checking~\citep{choudhury:ddc}, we show here
  that it is sound to use \emph{any} observer level that is compatible with
  the terms being compared. Furthermore, this relation can be composed
  step-wise: if $[[a]]$ is equal to $[[b]]$ at level $[[psi1]]$ and $[[b]]$ is
  equal to $[[c]]$ at level $[[psi2]]$, then we can also equate $[[a]]$ and
  $[[c]]$, even though those terms may not be equal at $[[psi1 /\ psi2]]$.
  We define how
  dependency tracking works in \name{} in Section~\ref{sec:typing} and 
  its role in definitional equality in Section~\ref{sec:defeq}.

\item We internalize the indistinguishability judgment in \name{}
  as an indexed equality type (Section~\ref{sec:equality}). Programmers can
  use this type to reason about program equivalences at a specified observer
  level. For security type systems, the observer-indexed equality type can be
  used to show that two programs appear identical to observers with lower
  clearance levels, ensuring that no secret information is leaked after, say,
  a significant refactoring.  In a system with relevance tracking, programmers
  can reason about programs without concerning themselves with the irrelevant
  type-level computations that are meant to be erased after compilation.
  % \jws{Removed word to avoid the orphan.}
  %
  This design crucially depends on the flexible definition of equality above:
  the elimination rule for this type must be able to reason about
  convertibility at any observer level.

\item To ensure that our design makes sense, we have mechanically checked its
  metatheoretic properties, including type soundness and noninterference. Our
  type soundness result requires showing the consistency of definitional
  equality, or a proof that equal types have the same head form. We provide an
  overview of these proofs in Section~\ref{sec:metatheory}.

\item We identify constraints on dependency tracking that are imposed
  by our use of indistinguishability for type conversion and discuss these 
  constraints in Section~\ref{sec:discussion}. 
  Our design contrasts with prior
  work~\citep{abadi1999core,shikuma2008proving, choudhury:ddc}
  that we have found to be incompatible with our approach. The insights
  gained in exploring this design space also explain why several
  existing mechanisms for tracking run-time irrelevance in dependently typed
  languages cannot be extended to compile-time irrelevance.
  % \jws{Removed words to avoid the orphan.}

\item To understand how \name works in practice, we have developed a
  prototype implementation that extends our core language with level-indexed data
  types. We have used this implementation to experiment with the development
  of sample programs. In the next section, we use these examples to provide an
  introduction to \name. We discuss our implementation in more detail in
  Section~\ref{sec:implementation}.

\end{itemize}

Our Coq proof scripts and prototype implementation are available
as supplementary materials.
%\jws{Removed ``to reviewers'' which feels like a strange thing to put in the final version.}

\section{Motivating Examples}
\label{sec:examples}

In this section,
we present three applications of dependency analysis as motivating examples:
run-time irrelevance\punctstack{,}\footnote{\dotpi{impl}{pi/RunTime.pi}}
compile-time irrelevance\punctstack{,}\footnote{\dotpi{impl}{pi/CompileTime.pi}}
and secure information flow\punctstack{.}\footnote{\dotpi{impl}{pi/InfoFlow.pi}}

All examples in this section have been type checked by our implementation,
with footnotes pointing to the corresponding source file. For clarity, we omit
some type and level annotations and use the concrete syntax of \name{} as
presented in the next section.

To make these examples more realistic, we use data definitions that are not in
the \name{} core calculus but are supported by our prototype implementation
(Section~\ref{sec:implementation}).  We also use constructs desugared from
\name{}'s dependent pairs: nondependent pairs; dependent pairs where one
component is labeled with a different level; and box types $[[T psi A]]$,
similar to the graded modal type of \citet{abadi1999core}, that encapsulate
terms at a higher level.

\subsection{Run-Time Irrelevance}
\label{sec:runtimeirrel}

Parts of programs that are not computationally meaningful when run
can be considered \textit{run-time irrelevant}.
An optimizing compiler can erase the run-time irrelevant parts
to yield a more space- and time-efficient program,
because those arguments do not need to be kept around or reduced needlessly
~\citep{paulinmohring1989coc,brady:phd}.
Run-time irrelevance annotations tell the compiler what can be erased,
and in a well typed program, irrelevant terms are never used relevantly.

For instance, consider the (simplified) polymorphic constant function,
whose type argument and unused input argument we mark irrelevant with label $[[I]]$. Unlabeled
arguments and definitions are treated as relevant by default,
sometimes labeled with $[[R]]$ for clarity, where  $[[R < I]]$.
% For instance, consider the polymorphic identity function,
% whose type argument we mark irrelevant with label $[[I]]$.
% Unlabeled arguments and definitions are treated as relevant by default,
% sometimes labeled with $[[R]]$ for clarity, where $[[R < I]]$.
\begin{lstlisting}
    const :$^[[R]]$ (A :$^[[I]]$ Type) -> A -> A$^[[I]]$ -> A
    const = \A :$^[[I]]$ Type. \x : A. \y :$^[[I]]$ A. x
\end{lstlisting}
The $[[R]]$ annotation next to \lstinline|const| says that
\lstinline|const| is a relevant computation. Since \lstinline|const|
is a function, the relevance of \lstinline|const| is determined by the
relevance of its output. That is, the output of a relevant function
can only use an irrelevant argument in irrelevant positions, as part
of a type annotation (as is the case for the type argument
\lstinline|A|) or as an argument to another function that takes an
irrelevant argument. If we try to return the irrelevant argument
\lstinline|y| instead as the output for \lstinline|const|, the
program will no longer type check.
Since well-typedness enforces that the first and third arguments of
\lstinline|const| cannot be used relevantly in its body, both
arguments can be safely erased at run-time.
% The type argument \lstinline|A| and the variable
% \lstinline|y| of type \lstinline|A| are all annotated with
% $[[I]]$. The type argument \lstinline|A| appears syntactically in the
% body of \lstinline|const| as the type annotation of
% \lstinline|y|. 

% The signature reads that \lstinline!id! is a relevant computation that
% takes an irrelevant type argument \lstinline!A! and a relevant
% argument of type \lstinline|A| as its inputs 

% Because the type argument only appears in the body of the function as
% a type annotation for the argument \lstinline!x!,
% it can safely be erased.
% On the other hand, if we attempt to ascribe \lstinline!id! with the
% type \lstinline!(A :$^[[I]]$ Type) -> A$^[[I]]$ -> A!, the program
% will be rejected since it attempts to return an irrelevant input when
% \lstinline!id! is itself a relevant computation that does not have
% access to arguments at level $[[I]]$.

% accidentally label the second argument as irrelevant,
% well-typedness enforces that the entire function must be used irrelevantly as well,
% since it returns an irrelevant term.
% \begin{lstlisting}
%     id' :$^[[I]]$ (A :$^[[I]]$ Type) -> A$^[[I]]$ ->  A
%     id' = \A :$^[[I]]$ Type. \x :$^[[I]]$ A. x
% \end{lstlisting}

Run-time irrelevance is also useful for extracting only the relevant parts out of proof-carrying code.
Consider the following inductive type representing the proposition that a natural is even.
Its constructors are labeled as irrelevant because we intend to erase such proofs.

\begin{lstlisting}
    data Even : Nat -> Type where
      ZEven  :$^[[I]]$ Even 0
      SSEven :$^[[I]]$ (n : Nat) -> (Even n)$^[[I]]$ -> Even (Succ (Succ n))
\end{lstlisting}

To represent the type of even naturals,
we use the subset type \lstinline!{x :$^[[R]]$ Nat | Even x}!.
Its elements are relevant dependent pairs whose first component is a relevant natural,
and whose second component is an irrelevant proof of evenness.
Using a lemma stating that the sum of two even naturals is even (definition elided),

\begin{lstlisting}
    evenEven :$^[[I]]$ (n : Nat) -> (m : Nat) -> (Even n)$^[[I]]$ -> (Even m)$^[[I]]$ ->  Even (add n m)
\end{lstlisting}

we can define addition on even naturals.
If all its run-time irrelevant parts---%
the proofs of evenness---%
are erased, this function corresponds exactly to the usual addition on naturals.
% For concision below and elsewhere, we use the pattern-matching construct supported by our implementation.
\scw{why do we say that we are using pattern matching? shouldn't that be implied by datatypes?
And should we say that our implementation doesn't do type inference?}
\jc{this is pattern matching as in I'm writing the arguments of functions
to the left of the equals, which the implementation doesn't have.
I removed the comment because pattern-matching notation is so common
people probably aren't even going to notice they're reading it.}

\begin{lstlisting}
    addEven : {x : Nat | Even x} -> {x : Nat | Even x} -> {x : Nat | Even x}
    addEven (n, en) (m, em) = (add n m, evenEven n m en em)
\end{lstlisting}

Consider next the following inductive type for length-indexed vectors,
whose length index is labeled as irrelevant.
These vectors can be erased to ordinary inductive lists.

\begin{lstlisting}
    data Vector (A :$^[[I]]$ Type) :$^[[I]]$ Nat$^[[I]]$ ->  Type where
      Nil  : Vector A 0
      Cons : (n :$^[[I]]$ Nat) -> A -> Vector A n -> Vector A (Succ n)
\end{lstlisting}

If we were to implement a map function over vectors,
the index of the vector ensures that it returns a vector with the same length,
and by marking the index as irrelevant,
the function can be erased at run time to one that looks just like a map function on lists.

\begin{lstlisting}
    map : (A :$^[[I]]$ Type) -> (B :$^[[I]]$ Type) -> (n :$^[[I]]$ Nat) ->
          (A -> B) -> Vector A n -> Vector B n
    map A B n f v = case v of
      Nil => Nil
      Cons n' y ys => Cons n' (f y) (map A B n' f ys)
\end{lstlisting}

It is worth pointing out that the irrelevant arguments \lstinline|A|,\lstinline|B|, and
\lstinline|n| do appear syntactically in the body of
\lstinline|map|. However, this is allowed since they all appear as irrelevant
arguments to a function, which, in this case, is \lstinline|map|
itself.

\subsection{Compile-Time Irrelevance}

Relevance tracking is not only useful for run-time erasure, but also
for ignoring terms during type conversion, known as \emph{compile-time
irrelevance}~\citep{debruijn1994automath}. It allows us to
type check the following program where the second
argument is a type constructor that ignores its input.
%\jws{Removed word to avoid orphan.}
\begin{lstlisting}
    substNat : (n :$^[[I]]$ Nat) -> (P :$^[[Low]]$ Nat$^[[I]]$ ->  Type) -> P n -> P 0
    substNat n P e = e
\end{lstlisting}
In \lstinline{substNat}, the type checker converts the type of the
argument \lstinline{e} from \lstinline{P n} to the return type \lstinline{P 0}.
This is valid because \lstinline{P} is typed at level $[[Low]]$ with
an input typed at level $[[High]]$, so these terms are indistinguishable to an
$[[Low]]$ observer.

However, it is not always possible to type a type constructor at level $[[Low]]$.
Recall \lstinline{Vector}, which takes in a type parameter and length index that
we regard as run-time irrelevant, labeled $[[High]]$. If we tried to type the overall
function \lstinline{Vector} at level $[[Low]]$, then
\lstinline{Vector Bool 1} and \lstinline{Vector Nat 2} would be indistinguishable,
which would violate type safety since we can now extract an element
out of a boolean vector and treat it as a natural number.
To ensure we do not convert between these two types,
\lstinline{Vector} must be typed at level at least $[[High]]$.

With only two levels $[[Low]]$ and $[[High]]$,
we cannot describe what is irrelevant to a type
already typed at level $[[High]]$. Being able to talk about
irrelevance at level $[[High]]$ is useful in practice: we may want to treat
proofs of some propositions as compile-time irrelevant or erase
irrelevant components to speed up type-level reduction.
We can add a label $[[S]]$ (for super-high) where
$[[R]] < [[I]] < [[S]]$. This allows us to define the following
program where the evenness proof, typed at level $[[S]]$, is ignored
by a type constructor \lstinline{P} typed at level $[[I]]$.

\begin{lstlisting}
    substEven : (P :$^[[I]]$ Nat$^[[I]]$ ->  (Even n)$^[[S]]$ ->  Type) -> (n :$^[[I]]$ Nat) ->
                (en$_1$ :$^[[S]]$ Even n) -> (en$_2$ :$^[[S]]$ Even n) -> P n en$_1$ ->  P n en$_2$
    substEven n P en$_1$ en$_2$ p = p
\end{lstlisting}

Intuitively, with the ordering of relevance labels $[[R]] < [[I]] < [[S]]$,
terms at higher relevance labels are indistinguishable
from the perspective of observers at lower relevance labels.
Since the function \lstinline{substEven} is typed at level
$[[R]]$, all of its argument except for \lstinline{p} are
erasable at run time. The type constructor \lstinline{P},
on the other hand, is typed at
level $[[I]]$. In \lstinline{P n en$_1$} and
\lstinline{P n en$_2$}, the natural \lstinline{n} is typed at
level $[[I]]$ and used relevantly at compile time, whereas the evenness proofs
\lstinline{en$_1$} and \lstinline{en$_2$} are both typed at level
$[[S]]$ and used irrelevantly at compile time. As a result, we can convert between
evenness proofs, but the natural must remain the same.

\name{} also includes a propositional equality type $a \sim^[[psi]] b$ 
that internalizes indistinguishability observed at level $[[psi]]$.
The $[[refl]]$ term introduces elements of this type. For example, 
\lstinline{P en$_1$ ~$^[[I]]$ P en$_2$} is inhabited by $[[refl]]$,
since the two arguments are
indistinguishable to $[[I]]$ observers.

This level-indexed equality type enables reasoning about equalities at that level.
For example, suppose we have a lemma showing commutativity of natural numbers (definition elided).
\begin{lstlisting}
    addComm : (n : Nat) -> (m : Nat) -> add n m ~$^{[[R]]}$ add m n
\end{lstlisting}
We can use this lemma with $[[transp]]$, the elimination term for
propositional equality, to prove the commutativity of adding two even
naturals.

\begin{lstlisting}
    addEvenComm : (en : {n : Nat | Even n}) -> (em : {m : Nat | Even m}) ->
                   addEven en em ~$^{[[R]]}$ addEven em en
    addEvenComm (n, en) (m, em) = $[[transp]]$ (addComm n m) $[[refl]]$
\end{lstlisting}

Given \lstinline{(n, en)} and \lstinline{(m, em)},
our goal is to show that \lstinline{(add n m, evenEven n m en em)}
and \lstinline{(add m n, evenEven m n em en)} are propositionally equal.
In addition to showing that \lstinline{(add n m)} and \lstinline{(add m n)} are equal,
we would normally need to manually prove that \lstinline{(evenEven n m en em)}
and \lstinline{(evenEven m n em en)} are equal,
further complicated by their definitionally unequal types \lstinline{Even (add n m)}
and \lstinline{Even (add m n)}.
However, since these proofs are compile-time irrelevant,
and we wish to prove a relevantly-indexed equality,
the type checker ignores the irrelevant second components,
and it suffices to prove the relevant first components equal.

\subsection{Information Flow}

Dependency analysis can be used to model \textit{secure information flow}~\citep{denning1977infoflow},
where information can never flow from a higher security level $[[High]]$
to a lower security level $[[Low]]$.
To an observer at $[[Low]]$, values of a given type at $[[High]]$ are
indistinguishable from one another,
and higher-level information therefore cannot be used meaningfully
except to be passed around.
On the other hand, observers with higher-level clearance are able to use lower-level information.
In the following examples, unlabeled arguments and definitions are treated as low-security by default.

Consider the construct $[[box High a]]$
that hides a secret value $[[a]]$ of type $[[A]]$
in a high-security box of type $[[T High A]]$. 
One way a low-clearance function can inspect and branch on a secret natural,
for instance, is to always box up its output.

\begin{lstlisting}
    secretPred :$^[[Low]]$ (n :$^[[High]]$ Nat) -> $[[T High Nat]]$
    secretPred n = $\ottkw{box}^[[High]]$ (case n of
      Zero => Zero
      Succ m => m)
\end{lstlisting}

Unboxing a boxed secret value yields a secret value,
which a low-clearance function can't do much with except to pass it to a function
that takes a high security value, such as in the \lstinline{ap} function below.
All of \lstinline{ap}'s arguments are low-security and it returns a low-security value,
but it can still run a secret function on a secret value
without ever knowing what was done and to which value.

\begin{lstlisting}
    ap :$^[[Low]]$ (A : Type) -> (B : Type) -> $T^[[High]]$ (A$^[[High]]$ ->   B) -> $T^[[High]]$ A ->   $T^[[High]]$ B
    ap A B f a = $\ottkw{box}^[[High]]$ ($\ottkw{unbox}^[[High]]$ f ($\ottkw{unbox}^[[High]]$ a))
\end{lstlisting}

We can show metatheoretically that low-clearance functions cannot distinguish between high-level arguments,
using the noninterference property of \name{} (Section~\ref{sec:noninterference}).
Furthermore, with internalized indistinguishability, programmers can also reason about noninterference within the language.
% \scw{Maybe we should stop here and not show indist. Let readers use their imagination.}
% \jc{Exercise for the reader!}

\section{\longname}
\label{sec:typing}
\begin{figure}[h]
\[
\begin{array}{lcll}
% \mathit{Levels, sorts}\\
% [[psi]]     & ::= & \ldots  & \mbox{levels (a set $\mathcal{L}$ forming a lattice)} \\
% [[s]]       & ::= & \ldots  & \mbox{sorts (a set $\mathcal{S}$)} \\ \\
\mathit{Contexts}\\
[[W]]       & ::= & [[empty]]\ |\ [[W ++ x : psi A]] & \mbox{variables
                                                       annotated by
                                                       level and type} \\
[[P]]       & ::= & [[empty]]\ |\ [[P ++ x : psi]] & \mbox{variables
                                                       annotated by level} \\ \\
\mathit{Terms}\\
[[a]],[[b]],[[A]],[[B]] & ::= & [[s]]\ |\ [[x]]\  |\ [[Unit]] \ |\ [[unit]]
                  & \mbox{sort, variables, unit type, unit} \\
            & |   & [[Pi x : psi A. B]]\ |\ [[\psi x : A . a]]\ |\ [[a psi b]]
                  & \mbox{function types, abstractions, applications} \\
            & |   & [[x : psi A && B]]\ |\ [[(a psi & b)]]\ |\ [[fst  psi a]]\ |\ [[snd  psi a]]
                  & \mbox{dependent pair types, pairs, projections} \\
            & |   & [[psi a ~ b : psi0 A ]]\ |\  [[refl]]\ |\ [[transp
                    a b]]
                  & \mbox{equality types, reflexivity proof, transport} \\
            & |   & [[Bool]]\ |\  [[true]]\ |\  [[false]]\ |\  [[if a b0 b1]]
                  & \mbox{boolean type, true, false, if}
% \mathit{Values}\\
% [[v]] & ::= & [[s]]\ |\ [[Unit]] \ |\ [[unit]]
%                   & \mbox{sort, unit type, unit} \\
%             & |   & [[Pi x : psi A. B]]\ |\ [[\psi x : A . a]]
%                   & \mbox{function types, abstractions} \\
%             & |   & [[x : psi A && B]]\ |\ [[(a psi & b)]]
%                   & \mbox{dependent pair types, pairs} \\
%             & |   & [[psi a ~ b : psi0 A ]]\ |\  [[refl]]
%                   & \mbox{equality types, reflexivity proof}
\end{array}
\]
  \caption{Syntax of \name}
  \label{fig:syntax}
\end{figure}
\ifextended
\begin{figure}[t]
\begin{tabular}{lll}
\textbf{\emph{Statics}} & %\textit{(Sections \ref{sec:modes}, \ref{sec:relevance}, \ref{sec:dec} and \ref{sec:naturals}) } 
\\
$[[W |- psi a : A]]$ & \textit{Typing} & Figures~\ref{fig:typing},
                                         \ref{fig:typingsum}, \ref{fig:typingeq}  \\
$[[|- W]]$ & \textit{Context formation} & Figure~\ref{fig:typing}   \\
$[[P |- a == b]]$ & \textit{Definitional equality} & Figure~\ref{fig:defeq}  \\
$[[P |- psi a == b]]$ & \textit{Indistinguishability} & Figure~\ref{fig:defeq}  \\
$[[P |- psi psi0 a == b]]$ & \textit{Guarded indistinguishability} & Figure~\ref{fig:defeq}  \\
%\\
\\
\textbf{\emph{Dynamics}}& %\textit{(Section \ref{sec:semantics}) }
\\
$[[a ~> b]]$ & \textit{Reduction} & Figure~\ref{fig:red}   
\\
\end{tabular}
\caption{Summary of \name{} judgment forms}
\label{fig:judgments}
\end{figure}
\fi

\begin{figure}[h]
\begin{minipage}{0.9\textwidth}
\drules[T]{$[[W |- psi a : A]]$}{Typing}{Type, Conv, Var, Pi, Abs,
  App, TyUnit, TmUnit, Bool, True, False, If}
\drules[Ctx]{$[[ |- W]]$}{Context well-formedness}{Empty, Cons}
\end{minipage}
\caption{Typing rules (selected)}
\label{fig:typing}
\end{figure}

In this section, we present the syntax, judgments, and derivation rules of
\name{}.  The syntax of \name{} (Figure~\ref{fig:syntax}) is parameterized
over a lattice $\mathcal{L}$ of dependency levels (following \citet{abadi1999core}) and a
set of sorts $\mathcal{S}$ (following \citet{barendregt:lambda-calculi-with-types}).  We use the metavariable
$[[psi]]$ for dependency levels and $[[s]]$ for sorts. For lattices, we use
$[[psi1 /\ psi2]]$ and $[[psi1 \/ psi2]]$ for the meet and join operations.  For
the examples in this section, we instantiate the lattice to $\{[[Low]], [[High]]\}$, where
$[[Low]] < [[High]]$, and the set of sorts to include the constant $[[star]]$, which
represents the type of types.

\ifextended
The judgment forms for \name are summarized in Figure~\ref{fig:judgments}.
\fi
The typing judgment of \name takes the form
$[[W |- psi a : A]]$.
Ignoring the labels $[[psi]]$ in the judgment
form and its rules, the typing rules in
Figure~\ref{fig:typing} are a variant of Barendregt's Pure Type Systems (PTS)~\citep{barendregt:lambda-calculi-with-types}.
These rules are parameterized over a set of sorts $\mathcal{S}$, a set of axioms $\mathcal{A}\subseteq \mathcal{S}^2$,
and sets of rules $\mathcal{R}_\Pi \subseteq \mathcal{S}^3$,
$\mathcal{R}_U \subseteq \mathcal{S}$, $\mathcal{R}_B \subseteq \mathcal{S}$.
The axioms dictate how sorts are typed in \rref{T-Type},
and the rules dictate the sorts assigned to types.
For example, function types are typed with respect to
the sorts of their domain and codomain types in \rref{T-Pi}.
The sorts, axioms, and rules can determine the expressiveness
of the language and its normalization properties.
For example, by instantiating $\mathcal{S}$ to ${[[star]]}$, $\mathcal{A}$ to $\{([[star]],[[star]])\}$, and
$\mathcal{R}_\Pi$ to $\{([[star]],[[star]],[[star]])\}$, we obtain a
system with unrestricted access to type-level computation,
but which fails the normalization property due to Girard's paradox~\citep{girard-thesis}.\footnote{In Section~\ref{sec:normalization}, we discuss our
conjectures about the normalization behavior of various instantiations of \name{}.}

\Name supports dependency tracking, as described in Section \ref{sec:dep}. This leads
to the most significant divergence from Barendregt's system, the definitional equality
relation used in \rref{T-Conv}.
This rule permits converting the type of a term to a definitionally equal
type, but whereas Barendregt uses $\beta$-equivalence, \name{} makes use of the
indistinguishability relation introduced in Section~\ref{sec:defeq}.
We also extend the core system with dependent pairs
(Section~\ref{sec:pair}) and propositional equality
(Section~\ref{sec:equality}). The boolean type $[[Bool]]$ is used as a
base type to formulate the noninterference property (Section~\ref{sec:noninterference}).
For concision, we use $[[A psi -> B]]$ as a shorthand for $[[Pi x :
psi A . B]]$ when $[[x]]$ does not appear in $[[B]]$. 

\subsection{Dependency Tracking}\label{sec:dep}
The level $[[psi]]$ from the typing judgment $[[W |- psi a : A]]$ 
indicates the \emph{observer level} used to type check the term
$[[a]]$. As enforced in the precondition of \rref{T-Var}, to access a variable,
the observer level of the computation must be at least as high as the level
at which the variable is labeled. 

Functions that type check at a lower observer level may still accept a higher-level argument,
as long as the argument is only used in higher-level positions.
\Rref{T-Abs} constructs a function that expects a computation
with observer level $[[psi0]]$ as its argument. Correspondingly, \rref{T-App}
checks the function argument at level $[[psi0]]$, a level that is
completely independent of the observer level $[[psi]]$. Both the
abstraction and application forms are annotated with this argument level.

As an example of how \rref{T-Var, T-Abs, T-App}
support relevance tracking, consider the following term with a hole ($[[hole]]$).
\[ [[\High x0 : Bool . \Low x1 : (Bool High -> Bool) . hole]]   \]
If we want to type check the term at level $[[Low]]$, we cannot fill
$[[hole]]$ with $[[x0]]$ since the observer level required to access
$[[x0]]$ is $[[High]]$. However, if we substitute $[[x1 High x0]]$ for
$[[hole]]$, the term type checks at $[[Low]]$ despite $[[x0]]$
appearing in the function body; when checking $[[x0]]$ as an argument
to the function $[[x1]]$, \rref{T-App} changes the observer level from
$[[Low]]$ to $[[High]]$ so $[[x0]]$ can be legally used.
\Rref{T-Abs} ensures that all functions that type
check at level $[[Low]]$ with type
$[[Bool High -> Bool]]$ behave like constant functions. The function
$[[\High x : Bool . x]]$, for example, is ill-typed at level $[[Low]]$
since $[[x]]$ is at level $[[High]]$ whereas the observer level is
$[[Low]]$. In other words, the declaration $[[x1 : Low Bool High ->
Bool]]$ encodes the contract that $[[x1]]$ does not use its input to
construct its result in a way that can be observed at level
$[[Low]]$. This gives an intuitive justification for why the term $[[\High
x0 : Bool . \Low x1 : (Bool High -> Bool) . x1 High x0]]$, which type
checks at $[[Low]]$, does not leak the high-level information from $[[x0]]$.

The observer level $[[psi]]$ at the colon does not have to be
precise. The following subsumption lemma says if a term type checks at
level $[[psi]]$, then it also type checks at levels higher than $[[psi]]$.
It can be viewed as a propagation of variable level subsumption $[[psi0 <= psi]]$
from \rref{T-Var} to the rest of the system.
In terms of relevance tracking,
this lemma embodies the idea that relevant computations can be lifted
to become irrelevant.
\begin{lemma}[Subsumption\footnote{\dotv{proofs/upgrade.v}{upgrade.v}{typing\_subsumption}}]
  \label{lemma:subsumption}
  If $[[W |- psi0 a : A]]$
  and $[[ psi0 <= psi ]]$,
  then $[[W |- psi a : A]]$.
\end{lemma}

In \rref{T-Abs}, we check that the function type can be assigned a sort
$[[s]]$. The most interesting part of this premise
is the choice of the function type level $[[psi1]]$, which is independent of
the observer level $[[psi]]$ and the argument level $[[psi0]]$. To show
that a type is well formed, it suffices to show that the type is
well formed at some arbitrarily chosen level.
The flexible choice of the
observer level for the well-formedness check affects how we formulate
the regularity property.

\begin{lemma}[Regularity\footnote{\dotv{proofs/regularity.v}{regularity.v}{typing\_regularity}}]
  \label{lemma:regularity}
  If $[[W |- psi a : A]]$,
  then there exists a level $[[psi0]]$ and a sort $[[s]]$ such that
  $[[W |- psi0 A : type s]]$.
\end{lemma}
If we instantiate $[[psi]]$ to a bounded lattice, Lemma~\ref{lemma:regularity} is
equivalent to the statement that the types of terms are always well typed at
the topmost level.
Intuitively, we are allowed to check well-formedness
at our level of choice because given $[[W |- psi a : A]]$, \name{}
does not have an operator to allow the retrieval of type $[[A]]$
from the term $[[a]]$ at run time.
As a result, we can safely allow $[[A]]$ to rely on secrets that are
not accessible to $[[a]]$ without worrying about any leakage.

\subsection{Dependent Pairs}\label{sec:pair}

\begin{figure}[h]
\begin{minipage}{0.9\textwidth}
  \begin{tabular}{c}
    \fbox{$[[W |- psi a : A]]$}\hfill\textit{(Typing: dependent pairs)}
    \\ \\
    \drule[width=2.5in]{T-SSigma} \qquad \drule[width=2.5in]{T-SPair} \\ \\
    \drule[width=2.5in]{T-ProjOne} \qquad \drule[width=2.5in]{T-ProjTwo}
  \end{tabular}
\end{minipage}
\caption{Typing rules for dependent pairs}
\label{fig:typingsum}
\end{figure}

The typing rules for dependent pairs can be found in
Figure~\ref{fig:typingsum}. Dependent pair types take the form $[[x : psi0 A && B]]$.
To govern the formation of these types,
we extend our rule sets to include
$\mathcal{R}_\Sigma$, the set of rules dictating which pair types are valid,
analogous to $\mathcal{R}_\Pi$.
The label $[[psi0]]$ in a type $[[x : psi0 A && B]]$ specifies the level
of the first component: in \rref{T-SPair}, the first component $[[a]]$
is checked at the labeled level $[[psi0]]$ rather than the observer
level $[[psi]]$. This is analogous to \rref{T-App} where the argument
of the function is also checked independently at the labeled
level.

Dependent pairs are eliminated by two
projection operators $[[fst psi0 a]]$ and $[[snd psi0 a]]$. \Rref{T-Proj1}
projects out the first component when the observer level $[[psi]]$ is
greater than or equal to the labeled level
$[[psi0]]$, while \rref{T-Proj2} projects out the second component.
The precondition
$[[W |- psi0 fst psi0 a : A]]$ in \rref{T-Proj2} ensures the
well-formedness of the type $[[B {fst   psi0 a / x}]]$.
This condition is admissible when
$[[psi <= psi0]]$. Given $[[W |- psi a : x : psi0 A && B]]$, we can derive
$[[W |- psi0 a : x : psi0 A && B]]$ by subsumption
(Lemma~\ref{lemma:subsumption}). We can then
apply \rref{T-Proj1} to obtain $[[W |- psi0 fst psi0 a :
A]]$. However, the condition is not admissible in the general case.
When $[[not (psi <= psi0)]]$, the pair $[[a]]$ may
depend on variables that are accessible at level $[[psi]]$ but not at
$[[psi0]]$, so the term $[[fst psi0 a]]$ is not necessarily
typable at level $[[psi0]]$ given $[[W |- psi a : x : psi0 A && B]]$
and $[[not (psi <= psi0)]]$.

% Note that this precondition does
%  not follow from the other precondition $[[W |- psi a : x :psi0 A && B]]$ because $[[a]]$ could be a computation that relies
% on variables that are accessible at level $[[psi]]$ but not at
% $[[psi0]]$.

The box type $T^{[[psi0]]}$ from Section~\ref{sec:examples} can be defined in terms of
a dependent pair type whose first component is the type being boxed,
and whose second component is unit.
%
\[
    [[T psi0 A]]  :=  [[x : psi0 A && Unit]] \qquad\qquad
    [[box psi0 a]] := [[(a psi0 & unit)]] \qquad\qquad
    [[unbox psi0 a]] := [[fst psi0 a]]
\]
%
\Rref{T-Proj1} enforces that we can only unbox a term if the observer level
$[[psi]]$ is greater than or equal to the box level $[[psi0]]$.
We can derive the following
admissible rules for boxes from the rules for dependent pairs,
where the relation $\mathcal{R}_T$ is defined in terms of $\mathcal{R}_\Sigma$.
%
\[ \drule[width=1.8in]{T-T} \  \drule[width=1.8in]{T-Box} \
  \drule[width=2in]{T-Unbox}\]
\[([[s1]],[[s2]]) \in \mathcal{R}_T \iff \exists [[s3]],
  ([[s1]],[[s3]],[[s2]]) \in \mathcal{R}_\Sigma \land s3 \in \mathcal{R}_U \]
%
\Rref{T-Box, T-Unbox} behave similarly to \rref{T-SPair, T-Proj1}
. \Rref{T-Box} allows us to place the term $[[a]]$
inside a box with level $[[psi0]]$ protection provided that $[[a]]$
can type check at level $[[psi0]]$. \Rref{T-Unbox} allows us to
release the data stored in the box, but requires that the observer
level is greater than or equal to the level associated with the box type.

The box type is a graded modal type. It was first used for dependency
tracking in DCC~\citep{abadi1999core} as an indexed monad type. Our
formulation of the box type is most similar to a later variation of
DCC called the Sealing
Calculus~\citep{shikuma2008proving}, which releases the boxed content
through an $\ottkw{unbox}$ operator rather than the monadic
$\ottkw{bind}$ operator from DCC.

\subsection{Dependency-Aware Definitional Equality}\label{sec:defeq}

\begin{figure}[h]
\begin{minipage}{0.9\textwidth}
\fbox{$[[P |- a == b]]$} \hfill \emph{(Definitional equality)}\\
\[ \drule{E-One} \qquad \drule[width=3in]{E-Trans} \]

\drules[GE]{$ [[P |- psi a == b]] $}{Indistinguishability}{Var, Type, Sym, Trans, App, Abs, AppAbs, Pi, TyUnit, TmUnit, ProjOneCong, ProjTwoCong, ProjOneBeta, ProjTwoBeta, SPair, SSigma}
\drules[CGE]{$ [[P |- psi psi0 a == b]] $}{Guarded indistinguishability}{Leq,NLeq}
\end{minipage}
\caption{Equality rules (selected)}
\label{fig:defeq}
\end{figure}

\Rref{T-Conv} allows us to convert between types that are
definitionally equal. 
The definitional equality judgment in Figure~\ref{fig:defeq} takes the form $[[P |- a == b]]$
where $[[P]]$ is an erased context that maps variables to their
levels. In \rref{T-Conv}, the erased context $[[|W|]]$ is obtained
from the typing context $[[W]]$ by dropping the type annotations.

Unlike PTS, definitional equality in \name{} is defined in terms of the
indistinguishability judgment, which takes the form $[[P |- psi a == b]]$.
Indistinguishability is a coarser relation than $\beta$-equivalence as it
takes advantage of the dependency information to identify terms that
are not convertible through $\beta$-equivalence alone. For example, one
can derive the judgment $[[empty |- Low box High false == box High
true]]$. Even though the boxes store different boolean values, from an $[[Low]]$
observer's point of view, the two terms cannot be distinguished. Using
\rref{E-One}, we can inject indistinguishability into definitional
equality and show that $[[box High false]]$ and $[[box High true]]$ are
definitionally equal. 

\Rref{E-Trans} means that two terms are definitionally equal if they can be related by a
sequence of indistinguishability judgments, each of which may have a
distinct observer level. Note that there may not be a single observer level that relates all terms in the sequence.
In other words, for a fixed $[[P]]$, definitional equality can be defined as the transitive
closure of the relation $ R([[a]],[[b]])=\exists [[psi]], [[P |- psi a == b]]$. 
\scw{Can delete last sentence if we need space}

Indistinguishability $[[P |- psi a == b]]$ and guarded indistinguishability
$[[P |- psi psi0 a == b]]$ are mutually defined. 
In \rref{GE-App},
we use guarded indistinguishability to skip the comparison of
unobservable terms.
In particular, when using guarded indistinguishability to compare terms $[[a]]$ and $[[b]]$ that
are labeled at level $[[psi0]]$ but being observed at level $[[psi]]$, two cases
arise. Either $[[psi0]]$ is observable at $[[psi]]$, i.e. $[[psi0 <= psi]]$,
in which case we revert to indistinguishability due to \rref{CGE-Leq}, or $[[psi0]]$ is not
observable at $[[psi]]$, in which case we immediately conclude that the terms
are indistinguishable due to \rref{CGE-NLeq}. As a result, the lower the
observer level, the coarser, or more general, indistinguishability becomes, since there are more
opportunities where \rref{CGE-NLeq} is applicable.

\Rref{GE-AppAbs} is the $\beta$-rules for functions,
and the remaining rules make indistinguishability reflexive, symmetric, transitive, and congruent.
The box type, defined earlier as syntactic sugar for dependent pairs,
has the following admissible equality rules.
The example
$[[empty |- Low box High false == box High true]]$ from above
holds by \rref{GE-Box,CGE-NLeq} since $[[not (High <= Low)]]$.
%
\[ \drule[width=1.8in]{GE-T} \ \drule[width=1.8in]{GE-Box} \ \drule[width=2in]{GE-Unbox}
% \ \drule[width=1.8in]{GE-UnboxBox}
\]

Given $[[P |- High a == b]]$,
it is not necessarily the case that $[[P |- Low a == b]]$.
If $[[P |- psi a == b]]$, neither $[[a]]$ nor $[[b]]$ may use variables
that are unobservable at level $[[psi]]$. This restriction is
imposed by \rref{GE-Var} and rules related to projecting from pairs, which mirror similar 
constraints in the typing judgments.
When lowering $[[psi]]$, the typing might become invalid if the
observer level becomes too low to access a variable or project
out a dependent pair.
On the other hand, 
given $[[P |- Low a == b]]$, it is not necessarily the case that 
$[[P |- High a == b]]$. By raising $[[psi]]$, the equivalence becomes
finer, as more subterms are observable.
As a result,
given an arbitrary derivation of $[[P |- psi a == b]]$, we can neither
raise nor lower $[[psi]]$. 

Hence the label $[[psi]]$ in indistinguishability not only controls
how fine-grained the equality is but also records the usage of variables
in the two terms being equated, enforcing an untyped version of
dependency tracking. The coupling of the two purposes
within the same label may appear restrictive as we cannot easily
raise or lower the labels, but it is necessary to ensure that
\name{} is type sound. If we remove the restriction imposed by
\rref{GE-Var}, we could end up converting between two arbitrary types
$[[A]]$ and $[[B]]$ through this chain of equalities:
\[ [[A]] [[==]]^[[Low]] [[(\High x : type star . x) High A]] [[==]]^[[Low]] [[(\High x : type star . x) High B]] [[==]]^[[High]] B \]
Without the inequality constraint in \rref{GE-Var}, the
lambda term $[[\High x : type star . x]]$ would be allowed to appear in the first
equality labeled at $[[Low]]$. With the observer level set to
$[[Low]]$, we can derive the second equality without having to show
that $[[A]]$ and $[[B]]$ are equal. By instantiating $[[A]]$ and
$[[B]]$ to two types with distinct head forms, we could easily derive a
program that crashes in \name.

\subsection{Indexed Equality Type}\label{sec:equality}

\begin{figure}[h]
\begin{minipage}{\textwidth}
  \begin{tabular}{c}
    \fbox{$[[W |- psi a : A]]$}\hfill\textit{(Typing: equalities)}
    \\ \\
    \drule[width=1.5in]{T-CEq} \ \drule{T-Refl} \  \drule[width=2in]{T-Transp}
  \end{tabular}
\end{minipage}
\caption{Typing rules for indexed equalities}
\label{fig:typingeq}
\end{figure}

The indexed equality type internalizes the indistinguishability judgment
so programmers can reason about indistinguishability within \name itself.
The typing and equality rules related to
the equality type can be found in Figure~\ref{fig:typingeq} and \ref{fig:gdefeqeq}.
Ignoring the levels, the rules for the indexed equality rules in
\name{} closely correspond to the rules for identity types from
MLTT~\citep{Martin-Lof-1973}.

An indexed equality type takes the
form $[[psi0 a ~ b : psi1 A]]$ and is well formed when it satisfies
the conditions specified in \rref{T-CEq}. The binary relation
$\mathcal{R}_\sim$ is analogous to $\mathcal{R}_\Pi$ for function
types and determines the sort we can assign to an equality type.
Both $[[a]]$ and $[[b]]$
must have type $[[A]]$ with $[[psi1]]$ as their observer
level. 
The level $[[psi0]]$ indicates the observer level at which
$[[a]]$ and $[[b]]$ are compared and can be distinct from
$[[psi1]]$, though it must satisfy the constraint that $[[psi1 <=
psi0]]$. The $[[psi1 <= psi0]]$ constraint will be explained
when we introduce the elimination form for the equality type.
We defer the discussion of the second constraint $[[psi0 <=
psi]]$ to Section~\ref{sec:discussion} since its explanation requires
some background on the metatheoretic properties of \name{}.
Similar to the introduction rule of the identity type
from MLTT, a trivial $[[refl]]$ proof witnesses the reflexive
indistinguishability through \rref{T-Refl}, provided that the equality type is well formed.

\Rref{T-Transp} allows us to eliminate an equality proof based on the
(specialized) congruence property of indistinguishability:
for $[[psi0 <= psi]]$, if $[[P ++ x : psi0 |- psi B0 == B1]]$ and $[[P |-
psi a0 == a1]]$, then $[[P |- psi B0 {a0 / x} == B1 {a1 / x}]]$.
Given a motive
$[[B]]$ % containing a free variable $[[x]]$
and an equality proof $[[a]]$
between the terms $[[a0]]$ and $[[a1]]$, \rref{T-Transp} can convert a
term $[[b]]$ of type $[[B {a0 / x}]]$ to a term of type $[[B {a1 /
  x}]]$. The observer level $[[psi0]]$ from the equality type matches
the observer level
of the motive $[[B]]$, whereas the level $[[psi1]]$ matches the
level associated with the variable $[[x]]$.

The congruence property of indistinguishability
requires us carefully to constrain the observer levels so we do not
accidentally convert between incompatible types. If $[[B0]]$ and
$[[B1]]$ are indistinguishable at level $[[psi]]$, then we cannot
substitute in $[[a0]]$ and $[[a1]]$ if they are indistinguishable at
level $[[psi0]]$ such that $[[psi0]] < [[psi]]$. To see why this is
problematic, consider the following two judgments.
\[ [[x : Low  |- High unbox High x == unbox High x ]] \qquad
   [[empty |- Low box High false  == box High true ]] \]
The first judgment is true because both sides are identical. The
second judgment is true since $[[false]]$ and $[[true]]$ both appear in
high-level boxes and are thus indistinguishable for a low-level
observer.
If we were allowed to substitute the second equation for $x$ in the first
one, we would be able to derive $[[empty |- High unbox High (box High false) == unbox
High (box High true) ]]$. By reducing each side, we end up proving that
$[[empty |- High false == true]]$.
In general, if $[[B0]]$ and $[[B1]]$ are indistinguishable at level $[[psi]]$,
then we want to only substitute in an indistinguishability judgment that
is not too coarse for the level $[[psi]]$.

Finally, given an equality type $[[psi0 a ~ b : psi1 A]]$, due to the
constraint $[[psi1 <= psi0]]$ in \rref{T-CEq}, \rref{T-Transp} can only be
applied when $[[psi1 <= psi0]]$, so we are only allowed to eliminate an
indexed equality if the level $[[psi1]]$ of the term we are substituting is
bounded above by the level of the motive. This constraint does not limit the
expressiveness of \name{} but is convenient for our proofs. When
$[[psi1]] \nleq [[psi0]]$, the variable $[[x]]$ can only appear in $[[B]]$ in
contexts, such as the argument to a function, where the level is high enough
to use $[[x]]$. In that case, $[[B {a0 / x}]]$ and $[[B {a1 / x}]]$ are
indistinguishable at level $[[psi0]]$ by equating the subterms containing
$[[a0]]$ and $[[a1]]$ through \rref{CGE-NLeq} from guarded
indistinguishability. We can then assign type
$[[B {a1 / x}]]$ to $a$ by \rref{T-Conv} directly without using \rref{T-Transp}.  We
make this intuition formal in Section~\ref{sec:metatheory} through the full
congruence property of indistinguishability (Lemma~\ref{lemma:gdefeqcong}).
% \jws{This makes sense to me as an explanation for
%  why the restriction is not a problem, but I'm not yet sure why it's
%  necessary.}  \yl{That's correct. It doesn't explain the necessity. Necessity
%  can be explained in terms of Lemma~\ref{lemma:typinggdefeq} but it's way too
%  technical. I wonder if we can justify the necessity by saying that any
%  violations of the restriction unnecessarily burns the CPU cycles..}\scw{I said that 
%  it is convenient for our proofs, without being specific as to why.}

\begin{figure}[h]
\begin{minipage}{0.9\textwidth}
  \begin{tabular}{c}
    \fbox{$[[P |- psi a == b]]$}\hfill\textit{(Indistinguishability: equalities)}
    \\ \\
    \drule[width=4.0in]{GE-CEq} \ \drule[width=3.0in]{GE-TranspRefl} \\ \\
    \drule[width=3.0in]{GE-Transp} \ \drule[width=3.0in]{GE-Refl}
  \end{tabular}
\end{minipage}
\caption{Equality rules for indexed equalities}
\label{fig:gdefeqeq}
\end{figure}
Figure~\ref{fig:gdefeqeq} shows the indistinguishability rules
for indexed equality. When used in conjunction with \rref{GE-CEq},
\rref{T-Refl} allows us to use the $[[refl]]$ constructor to witness
the equality between the two terms $[[a]]$ and $[[b]]$ at observer level
$[[psi]]$ if we can show that $[[P |- psi a == b]]$. This is analogous
to how the reflexivity proof term can witness the equality between 
two $\beta$-equivalent terms in Martin-L{\"o}f type theory.
In \rref{GE-CEq}, the precondition requires that the observer level $[[psi0]]$ of
the equality type be lower than the observer level $[[psi]]$ of
the judgment. Moreover, the observer level used for the equalities in
the preconditions is $[[psi]]$ rather than $[[psi0]]$ even though
$[[psi0]]$ is the one labeled on the equality types. We will explain
the subtleties behind the choice of labels in \rref{GE-CEq} in
Section~\ref{sec:orderconstraints}.

\section{Metatheory}
\label{sec:metatheory}
\begin{figure}[h]
  \begin{minipage}{\textwidth}
    \drules[R]{$[[a ~> b]]$}{Reduction}{AppAbs, TranspRefl,
      ProjOneBeta, ProjTwoBeta, IfTrue, IfFalse}

    
    \[ [[v]]  ::= [[s]]\ |\ [[Unit]] \ |\ [[unit]] \ |\ [[Pi x : psi
      A. B]]\ |\ [[\psi x : A . a]] \ |\ [[x : psi A && B]]\ |\ [[(a
      psi & b)]] \ |\ [[psi a ~ b : psi0 A ]]\ |\  [[refl]]\ |\  [[Bool]]\ |\  [[true]]\ |\  [[false]] \]
  \end{minipage}
\caption{Reduction ($\beta$-rules only) and values}
\label{fig:red}
\end{figure}
In this section, we present our main results, including type soundness
(Theorem~\ref{theorem:preservation} and \ref{theorem:progress})
and noninterference (Theorem~\ref{theorem:noninterfere}),
with pointers to the corresponding Coq proof file and name in the footnote.
To formulate these properties, we first define the operational
semantics and the syntax of values in Figure~\ref{fig:red}. This figure
includes the $\beta$-rules only---the compatibility rules
are standard for a call-by-name language. The full reduction rules can
be found in the supplementary materials\footnote{\href{spec.pdf}{spec.pdf}}.

\subsection{Structural Lemmas}
\label{sec:structural}
The system supports the standard structural rules such as weakening
and narrowing. Since \rref{T-Conv} depends on the definitional
equality judgment, we need to first establish the structural lemmas for
definitional equality and indistinguishability before we can
derive the structural lemmas for the typing judgment.

\begin{lemma}[Indistinguishability Weakening\footnote{\dotv{proofs/defeq_weak.v}{defeq\_weak.v}{gdefeq\_weak\_nil}}]
  If $[[P |- psi a == b]]$,
  then $[[P ++ P0 |- psi a == b]]$.
\end{lemma}

\begin{lemma}[DefEq weakening\footnote{\dotv{proofs/defeq_weak.v}{defeq\_weak.v}{defeq\_weak\_nil}}]
  If $[[P |- a == b]]$,
  then $[[P ++ P0 |- a == b]]$.
\end{lemma}
In both weakening lemmas, we use $[[P ++ P0]]$ to indicate the
concatenation of two erased contexts with disjoint and unique domains.

The narrowing property says that equalities still hold
after lowering the levels of the variables in the context.
\begin{lemma}[Indistinguishability Narrowing\footnote{\dotv{proofs/narrow.v}{narrow.v}{gdefeq\_narrow}}]
  If $[[P |- psi A == B]]$ and $[[P0 <= P]]$, then
  $[[P0 |- psi A == B]]$.
\end{lemma}

\begin{lemma}[DefEq Narrowing\footnote{\dotv{narrow.v}{narrow.v}{defeq\_narrow}}]
  If $[[P |- A == B]]$ and $[[P0 <= P]]$, then
  $[[P0 |- A == B]]$.
\end{lemma}
The relation $[[P0 <= P]]$ compares pointwise the levels of two
erased contexts with the same domain.

Unlike for the weakening and narrowing properties, indistinguishability
and definitional equality behave differently when it comes to
congruence. The congruence lemma for indistinguishability is defined
as follows.
\begin{lemma}[GDefEq Congruence\footnote{\dotv{proofs/defeq_subst.v}{defeq\_subst.v}{gdefeq\_cong}}]
  \label{lemma:gdefeqcong}
  If $[[P ++ x : psi0 |- psi a0 == a1]]$ and $[[P |- psi psi0 b0 ==
  b1]]$, then $[[P |- psi a0 {b0 / x} == a1 {b1 / x}]]$.
\end{lemma}

It is crucial that the label $[[psi]]$ from the guarded
indistinguishability judgment matches the observer level from the
indistinguishability judgment. The guarded indistinguishability we are substituting in
must be as fine as the level of the indistinguishability judgment or
we may equate two terms with differing head forms. Given $[[P ++ x :
psi0 |- a0 == a1]]$, its derivation may consist of a sequence of
indistinguishability judgments with different observer levels. As a
result, we cannot formulate a congruence property for definitional
equality in a manner similar to Lemma~\ref{lemma:gdefeqcong}.
\begin{lemma}[DefEq Congruence\footnote{\dotv{proofs/defeq_subst.v}{defeq\_subst.v}{defeq\_cong}}]
  \label{lemma:defeqcong}
    If $[[P ++ x : psi0 |- a0 == a1]]$,  $[[P |- psi0 b0 ==
  b1]]$, and $[[b0]]$ is $\beta$-equivalent to $[[b1]]$, then $[[P |- a0 {b0 / x} == a1
  {b1 / x}]]$.
\end{lemma}
Since we do not know which observer levels are used underneath the derivation of
$[[P ++ x : psi0 |- a0 == a1]]$, we can only make the worst-case
assumption that we cannot skip any comparison between the subterms in $[[b0]]$ and
$[[b1]]$. That is why Lemma~\ref{lemma:defeqcong} additionally
requires that $[[b0]]$ and $[[b1]]$ are $\beta$-equivalent.

Indistinguishability does not include reflexivity as an
axiom, but we can prove that reflexivity holds for well-typed terms.
\begin{lemma}[Typing Indistinguishability\footnote{\dotv{proofs/typing_defeq.v}{typing\_defeq.v}{typing\_gdefeq}}]
  \label{lemma:typinggdefeq}
  % If $[[W |- psi a : A]]$, then for all $[[psi1]]$ such that $[[psi <= psi1]]$, 
  % we have $[[|W| |- psi1 a == a]]$.
  If $[[W |- psi a : A]]$, then $[[|W| |- psi a == a]]$. % Furthermore,
  % the derivation of $[[|W| |- psi a == a]]$ does not use \rref{CGE-NLeq}.
\end{lemma}

Now we are ready to prove some properties about the typing
judgments. The proofs follow the same structure as their
equality counterparts by induction over the derivation, though they
rely on the structural rules about equalities proven earlier.
% TODO: move narrowing and subsumption to the relevance tracking section
\begin{lemma}[Narrowing\footnote{\dotv{proofs/narrow.v}{narrow.v}{typing\_narrow}}]
  \label{lemma:narrow}
  If $[[W |- psi a : A]]$
  and $[[W0 <= W]]$,
  then $[[W0 |- psi a : A]]$.
\end{lemma}

Subsumption (Lemma~\ref{lemma:subsumption}) and weakening both follow
from Lemma~\ref{lemma:narrow}.
\begin{lemma}[Weakening\footnote{\dotv{proofs/weak.v}{weak.v}{typing\_weak\_nil}}]
  If $[[W1 |- psi a : A]]$
  and $[[|- W1 ++ W2]]$,
  then $[[W1 ++ W2 |- psi a : A]]$.
\end{lemma}

We can now prove the substitution property for the typing judgment
through structural induction over derivations.
\begin{lemma}[Substitution\footnote{\dotv{proofs/subst.v}{subst.v}{typing\_subst\_nil}}]
  \label{lemma:subst}
  If $[[W ++ x : psi0 A |- psi b : B]]$
  and $[[W |- psi0 a : A]]$,
  then $[[W |- psi b { a / x} : B {a / x} ]]$.
\end{lemma}
The \rref{T-Var} case of Lemma~\ref{lemma:subst} relies on
Lemma~\ref{lemma:subsumption} since we may substitute in a term with a
lower level when $[[psi0 <= psi]]$. The \rref{T-Conv} case
requires us to show that $[[ |W| |- A {a / x} == B {a / x} ]]$ given
$[[|W| ++ x : psi0 |- A == B]]$ and $[[W |- psi0 a : A]]$. This can be
done by composing Lemma~\ref{lemma:defeqcong} and
Lemma~\ref{lemma:typinggdefeq}.

\subsection{Preservation}
\label{sec:preservation}
\begin{figure}[t]
  \begin{minipage}{0.9\textwidth}
    \begin{tabular}{llll}
      \textbf{\emph{Parallel reductions}} & \\
      $[[P |- psi a => b]]$ & \textit{Indexed parallel reduction} & $\mapsto$
                                                                  & $[[P |- psi a == b]]$ \\
      $[[P |- psi psi0 a => b]]$ & \textit{Guarded parallel reduction} & $\mapsto$
                                                                  & $[[P |- psi psi0 a == b]]$ \\
      $[[P |- a => b]]$ & \textit{Parallel reduction} & $\mapsto$ 
                                                                  & $[[P |- a == b]]$
    \end{tabular}
  \end{minipage}
\caption{Summary of the parallel reduction judgments}
\label{fig:eqiffred}
\end{figure}
Similar to the preservation proof for PTS,
preservation in \name{} requires inversion lemmas about the typing judgment and
injectivity lemmas about the definitional equality judgment.

The inversion lemmas witness the fact that given $[[W |- psi v : A]]$, the
derivation must consist of the introduction rule for $[[v]]$ nested
under zero or more instances of the conversion rule. The proofs for
those lemmas, with the exception of the $[[refl]]$ form, can be
carried out by induction over the derivation. For now, we only show
the inversion lemma for abstractions and defer the explanation of the
$[[refl]]$ case until we discuss the injectivity properties; the rest
of the inversion lemmas follow the same pattern and are thus omitted.
\begin{lemma}[Inversion (Abs)\footnote{\dotv{proofs/inv.v}{inv.v}{typing\_abs\_inv}}]
  \label{lemma:absinv}
  If $[[W |- psi (\psi0 x : A . b) : A0]]$, then there exists a term
  $[[B]]$ such that the following conditions hold:
  \begin{itemize}
  \item $[[W ++ x : psi0 A |- psi b : B]]$
  \item $[[|W| |- Pi x : psi0 A . B == A0]]$
  \item $[[W |- psi1 A0 : type s]]$ for some level $[[psi1]]$ and some
    sort $[[s]]$
  \end{itemize}
\end{lemma}

To prove preservation for $\beta$-rules such as
\rref{R-AppAbs, R-TranspRefl} (Figure~\ref{fig:red}), we need to prove some injectivity lemmas about
definitional equality. For example, to prove that \rref{R-AppAbs} is
type preserving, we need to show that the two definitionally equal
function types must also have definitionally equal argument types and
return types. A similar proof can be found in the preservation proof
for PTS, where definitional equality is defined as
$\beta$-equivalence. 
To prove the injection lemma for function types
in a PTS, one needs to first prove the Church-Rosser
property, which states that two terms are $\beta$-equivalent if and only
if they can strongly reduce
to the same term within a finite number of
steps. The injection lemma then immediately follows since a
reduction sequence starting from a function type can only consist of
repeated reductions of its argument type or return type.

We adapt the proof technique and prove the equivalence between
the equality judgments and their corresponding parallel reduction
relations, summarized in Figure~\ref{fig:eqiffred}. Each
reduction relation corresponds to the equivalence relation in the
following sense: given two terms $[[a0]]$ and $[[a1]]$, the terms
$[[a0]]$ and $[[a1]]$ are equivalent if and only if there
exists some term $[[b]]$ such that $[[a0]]$ and $[[a1]]$ both reduce to
$[[b]]$ after a finite number of steps.
We omit the
rules for the parallel reduction relations since they closely correspond to the
equality rules in Figure~\ref{fig:defeq}. The rules for each
parallel reduction relation can be obtained from the
corresponding equivalence rule by removing the
transitivity and symmetry rules and replacing $\equiv$ with
$\Rightarrow$. These parallel reduction rules can be found in the
supplementary materials\footnote{\href{spec.pdf}{spec.pdf}}.
We also omit the structural lemmas for the reduction
relations since their specifications and proofs can be obtained with a similar replacement.


Similar to the indistinguishability judgment, given $[[P |- psi a =>
b]]$, neither $[[a]]$ nor $[[b]]$ is allowed to use variables from
$[[P]]$ that are not observable at level $[[psi]]$. The following
shows that if $[[P |- psi a => b]]$ or $[[P |- psi a == b]]$, then
indexed parallel reduction is reflexive on both $[[a]]$ and
$[[b]]$.

\begin{lemma}[Par Reflexive\footnote{\dotv{proofs/par.v}{par.v}{Par\_grade\_mutual}}]
  \label{lemma:pargrade}
  If $[[P |- psi a => b]]$, then $[[P |- psi a => a]]$ and $[[P |- psi
  b => b]]$.
\end{lemma}

\begin{lemma}[Indistinguishability Par Reflexive\footnote{\dotv{proofs/defeq\_proj.v}{defeq\_proj.v}{gdefeq\_grade}}]
  \label{lemma:gdefeqgrade}
  If $[[P |- psi a == b]]$, then $[[P |- psi a => a]]$ and $[[P |- psi
  b => b]]$.
\end{lemma}
Both lemmas follow from induction on the derivation of the premise.

Indexed parallel reduction generalizes $\beta$-reduction in the same way
that indistinguishability generalizes $\beta$-equivalence. It is possible
to embed the $\beta$-reduction relation into the parallel reduction
relation.
\begin{lemma}[Red Embed Par\footnote{\dotv{proofs/preservation.v}{preservation.v}{red\_embed}}]
  \label{lemma:redembed}
  If $[[a ~> b]]$ and there exists some $[[b0]]$ such that $[[P |- psi
  a => b0]]$, then $[[P |- psi a => b]]$.
\end{lemma}
In Lemma~\ref{lemma:redembed}, the precondition $[[P |- psi a => b0]]$
tells us that $[[a]]$ does not use variables with labels not
observable at $[[psi]]$. It provides us a context and a label at which
we can reduce from $[[a]]$ to $[[b]]$ through indexed parallel
reduction. The precondition $[[P |- psi a => b0]]$ can be supplied
through the well-typedness of $[[a]]$ by composing
Lemma~\ref{lemma:typinggdefeq} and Lemma~\ref{lemma:gdefeqgrade}.

From the structural lemmas, we can show that indexed parallel
reduction satisfies the following confluence property.
\begin{lemma}[Indexed Par Confluence\footnote{\dotv{proofs/par.v}{par.v}{Par\_confluence}}]
  \label{lemma:gparconf}
  If $[[P |- psi a => b0]]$ and $[[P |- psi a => b1]]$, then there
  exists some $[[b2]]$ such that $[[P |- psi b0 => b2]]$ and $[[P |-
  psi b1 => b2]]$.
\end{lemma}

Before we can use Lemma~\ref{lemma:gparconf} to derive the confluence
property of parallel reduction, we need to first show the following
downgrade property about indexed parallel reduction.
\begin{lemma}[Indexed Par Downgrade\footnote{\dotv{proofs/par.v}{par.v}{Par\_downgrade\_mutual}}]
  \label{lemma:gpardowngrade}
  If $[[P |- psi a => b0]]$ and $[[P |- psi0 a => b1]]$, then $[[P |-
  psi /\ psi0 a => b0]]$.
\end{lemma}
The term $[[b1]]$ does not appear in the conclusion. The only
information we need from the premise $[[P |- psi0 a => b1]]$ is that
$[[a]]$ has $[[psi0]]$ as its observer level. The proof proceeds by
induction over the derivation and leverages the fact that $[[psi /\
psi0]]$ is the greatest lower bound.

Confluence of parallel reduction then follows as a
corollary of Lemmas~\ref{lemma:gparconf} and \ref{lemma:gpardowngrade}.
\begin{lemma}[Par Confluent\footnote{\dotv{proofs/par.v}{par.v}{EPar\_confluence}}]
  \label{lemma:parconf}
  If $[[P |- a => b0]]$ and $[[P |- a => b1]]$, then there
  exists some $[[b2]]$ such that $[[P |- b0 => b2]]$ and $[[P |-
  b1 => b2]]$.
\end{lemma}
The proof proceeds by unfolding the definition of parallel
reduction. This gives us $[[P |- psi a => b0]]$ and $[[P |- psi0 a =>
b1]]$ for some $[[psi]]$ and $[[psi0]]$.  By applying
Lemma~\ref{lemma:gpardowngrade} twice, we can conclude that $[[P |- psi /\
psi0 a => b0]]$ and $[[P |- psi /\ psi0 a => b1]]$. The conclusion
then follows immediately by applying Lemma~\ref{lemma:gparconf}.

Once Lemma~\ref{lemma:parconf} is proven, we can easily show that
two terms are definitionally equal if and only if they parallel
reduce to the same term.
\begin{lemma}[DefEq Join\footnote{\dotv{proofs/defeq_proj.v}{defeq\_proj.v}{defeq\_ejoins\_iff}}]
  \label{lemma:defeqjoin}
$[[P |- a0 == a1]]$ if and only if there exists some $[[b]]$ such that $[[P
|- a0 =>* b]]$ and $[[P |- a1 =>* b]]$.
\end{lemma}
The backward direction is trivial since it is possible to embed
parallel reduction into definitional equality by simply unfolding definitions. The forward direction uses Lemma~\ref{lemma:parconf} to finish the
\rref{E-Trans} case and the rest of the cases are straightforward.

The injectivity lemmas for the function types and dependent sum types
follow as corollaries of Lemma~\ref{lemma:defeqjoin}.
\begin{lemma}[DefEq $\Pi$-Inj1\footnote{\dotv{proofs/defeq_proj.v}{defeq\_proj.v}{defeq\_pi\_proj1}}]
  \label{lemma:pifst}
  If $[[P |- Pi x : psi0 A0 . B0 == Pi x : psi0 A1 . B1]]$, then
  $[[P |- A0 == A1]]$.
\end{lemma}

\begin{lemma}[DefEq $\Pi$-Inj2\footnote{\dotv{proofs/defeq_proj.v}{defeq\_proj.v}{defeq\_pi\_proj2}}]
  \label{lemma:pisnd}
  If $[[P |- Pi x : psi0 A0 . B0 == Pi x : psi0 A1 . B1]]$, then
  $[[P ++ x : psi0 |- B0 == B1]]$.
\end{lemma}

\begin{lemma}[DefEq $\Sigma$-Inj1\footnote{\dotv{proofs/defeq_proj.v}{defeq\_proj.v}{defeq\_sigma\_proj}}]
  \label{lemma:sigmafst}
  If $[[P |- x : psi0 A0 && B0 == x : psi0 A1 && B1]]$, then
  $[[P |- A0 == A1]]$.
\end{lemma}

\begin{lemma}[DefEq $\Sigma$-Inj2\footnote{\dotv{proofs/defeq_proj.v}{defeq\_proj.v}{defeq\_sigma\_proj}}]
  \label{lemma:sigmasnd}
  If $[[P |- x : psi0 A0 && B0 == x : psi0 A1 && B1 ]]$, then
  $[[P ++ x : psi0 |- B0 == B1]]$.
\end{lemma}

Furthermore, Lemma~\ref{lemma:defeqjoin} allows us to show the
following strengthened inversion lemma about the indexed equality type.
\begin{lemma}[Refl GDefEq\footnote{\dotv{proofs/preservation.v}{preservation.v}{typing\_refl\_gdefeq}}]
  \label{lemma:reflgdefeq}
  If $[[W |- psi refl : psi0 a ~ b : psi1 A]]$, then $[[|W| |- psi0 a == b]]$.
\end{lemma}
In the proof of Lemma~\ref{lemma:reflgdefeq},
Lemma~\ref{lemma:defeqjoin} is used to show that given $[[|W| |- psi0
a0 ~ b0 : psi1 A == psi0 a ~ b : psi1 A ]]$, we can conclude that
$[[|W| |- psi0 a0 == a]]$ and $[[|W| |- psi0 b0 == b]]$.

Now, we can finally show that the reduction relation is type
preserving.
\begin{theorem}[Preservation\footnote{\dotv{proofs/preservation}{preservation.v}{preservation}}]
  \label{theorem:preservation}
  If $[[W |- psi a : A]]$
  and $[[a ~> b]]$,
  then $[[W |- psi b : A]]$.
\end{theorem}
The proof of Theorem~\ref{theorem:preservation} is carried out by
structural induction over the reduction relation. The inversion lemmas
are applied to the premise $[[W |- psi a : A]]$ to gain information
about the derivation. The injectivity lemmas are needed for the cases
that involve $\beta$-rules. The case for the second projection
needs some special treatment. As part of the proof, given $[[W |- psi a : x : psi0 A &&
B]]$ and $[[a ~> a0]]$, we need to prove that $[[snd psi0 a0]]$ can be
assigned the type $[[B {fst psi0 a / x}]]$. From \rref{T-Proj2} and
the inductive hypothesis, we can show that $[[W |- psi snd psi0 a0 : B {fst psi0 a0 / x}]]$.
To finish off the proof, we use Lemma~\ref{lemma:redembed} to inject
the reduction into the indistinguishability relation to show that
$[[B {fst psi0 a / x}]]$ is definitionally equal to $[[B {fst psi0 a0
  / x}]]$.

\subsection{Progress}
Before we prove progress, it is useful to derive the following simple
corollary of Lemma~\ref{lemma:defeqjoin}.
\begin{lemma}[DefEq Consistency\footnote{\dotv{proofs/defeq_par.v}{defeq\_par.v}{defeq\_consist}}]
  \label{lemma:defeqconsist}
 If $[[P |- v0 == v1]]$, then $[[v0]]$ and $[[v1]]$ have the same
 head form.
\end{lemma}
By Lemma~\ref{lemma:defeqjoin}, $[[v0]]$ and $[[v1]]$ can reduce to
the same term through parallel reduction. However, since parallel
reduction preserves head forms, $[[v0]]$ and $[[v1]]$ cannot have
conflicting head forms or it would be impossible to reduce them to the
same term.

Lemma~\ref{lemma:defeqconsist}, when combined with the inversion
lemmas from Section~\ref{sec:preservation}, can be used to prove the
following canonical form lemmas.
\begin{lemma}[$\Pi$-Canon\footnote{\dotv{proofs/progress.v}{progress.v}{typing\_pi\_canon}}]
  If $[[W |- psi v : Pi x : psi0 A . B]]$,
  then there exists some $[[A0]]$ and $[[a0]]$ such that
  $[[v]]=[[\psi0 x : A0 . a0]]$ .
\end{lemma}
\begin{lemma}[$\Sigma$-Canon\footnote{\dotv{proofs/progress.v}{progress.v}{typing\_ssigma\_canon}}]
  If $[[W |- psi v : x : psi0 A0 && B0]]$,
  then there exists some $[[a0]]$ and $[[a1]]$ such that $[[v = (a0 psi0
  & a1)]]$.
\end{lemma}
\begin{lemma}[Eq-Canon\footnote{\dotv{proofs/progress.v}{progress.v}{typing\_ceq\_canon}}]
  If $[[W |- psi v : psi0 b0 ~ b1 : psi1 A]]$,
  then $[[v]]=[[refl]]$.
\end{lemma}

We can now conclude our type soundness proof by showing the progress theorem.
\begin{theorem}[Progress\footnote{\dotv{proofs/progress.v}{progress.v}{progress}}]
  \label{theorem:progress}
If $[[W |- psi a : A]]$, then $[[a]]$ is a value or there
exists some $[[b]]$ such that $[[a ~> b]]$.
\end{theorem}
The proof proceeds by induction over the derivation $[[W |- psi a :
A]]$. Each non-value case can be discharged with the help of
its corresponding canonical form lemma.

Theorem~\ref{theorem:preservation} and \ref{theorem:progress} complete
our type soundness result. A simple corollary of the type soundness result is
the erasability of the levels that appear in the term syntax. For
example, the level $[[psi0]]$ in the application form $[[a psi0 b]]$
can block \rref{R-AppAbs} if it mismatches the level of
the lambda term. Preservation and progress imply that a
well-typed term's execution will never be blocked by mismatching
levels. Since blocking $\beta$-rules is the only way the levels
can affect reduction, we can safely erase the levels before running a
well-typed term. However, levels are still necessary for definitional
equality and indistinguishability because those relations are untyped. Refactoring 
to use a typed indistinguishability relation might allow us to remove the levels
from the term syntax completely, and we leave that as part of our future work.

\subsection{Noninterference}
\label{sec:noninterference}
In \name{}, one can formulate the following simulation property in terms of indistinguishability. It states
that programs that are indistugishable remain indistinguishable during evaluation.
\begin{lemma}[Simulation\footnote{\dotv{proofs/preservation.v}{preservation.v}{simulation}}]
  \label{lemma:simulation}
If $[[P |- psi a0 == b0]]$, $[[a0 ~>* a1]]$, and $[[b0 ~>* b1]]$, then
$[[P |- psi a1 == b1]]$.
\end{lemma}
Theorem~\ref{lemma:simulation} is easily derivable
from Lemma~\ref{lemma:redembed} and
Lemma~\ref{lemma:pargrade}. The conclusion of
Lemma~\ref{lemma:simulation} does not immediately tell us much about
the relation between $[[a1]]$ and $[[b1]]$ since indistinguishability
incorporates both reduction and irrelevance.

However, when combined with
Lemma~\ref{lemma:defeqjoin}, Lemma~\ref{lemma:simulation} imposes very strong
constraints on how two indistinguishable programs may behave after
evaluation.
With the boolean type, we formulate the
following noninterference theorem, which states that two indistinguishable
boolean computations reduce to the same boolean value.

\begin{theorem}[Noninterference\footnote{\dotv{proofs/progress.v}{progress.v}{non\_interference}}]
  \label{theorem:noninterfere}
If $[[ empty |- psi a == b ]]$, $[[empty |- psi a : Bool ]]$, $[[empty
|- psi b : Bool]]$, $[[a ~>* v0]]$, and $[[b ~>* v1]]$, then
$[[v0]] = [[v1]]$.
\end{theorem}
% Consider the case where $[[empty |- psi a0 == a1]]$, $[[a0 ~>* v0]]$, $[[a1 ~>* v1]]$,  $[[a0]]$ and $[[a1]]$ are both
% booleans. By composing Theorem~\ref{lemma:simulation} and the backward direction of
% Lemma~\ref{lemma:defeqjoin}, we deduce that $[[v0]]$ and $[[v1]]$ must
% both reduce to some common term, allowing us to conclude that
% $[[v0]]$ and $[[v1]]$ are syntactically equal.

% In the special case where $[[a0]]$ and $[[b0]]$ from
% Theorem~\ref{lemma:simulation} are both nautral numbers,
% Lemma~\ref{lemma:defeqconsist} ensures that both $[[a1]]$ and $[[b1]]$
% have the same head form.
% \yl{Is there anything else I need to say? The section is too short.}
% \scw{What about talking about how if $[[a]]$ is low at type $[[nat]]$ and mentions free high variable $[[x]]$, then $[[a { b1/ x } ]]$ and $[[a {b2 / x} ]]$ both evaluate 
% to the same number.}
\section{Discussion}
\label{sec:discussion}
\subsection{Incompatibility between Indistinguishability and the
  Upgrade Property}
\label{sec:upgrade}

Consider these variants of \rref{T-App, T-Abs},
with argument level $[[psi \/ psi0]]$ instead of $[[psi]]$.
\[ \drule[width=2.5in]{T-AppJoin} \qquad \drule[width=3in]{T-AbsJoin} \]
Compared to the corresponding rules from \name{},
the alternative rules are more permissive. The following judgment 
is derivable with the alternative rules but not the existing ones.
\[ [[ y : High Bool |- High (\Low x : Bool . x) Low y : Bool]]  \]
With \rref{T-App}, the application is illegal because the
argument $[[y]]$ is type checked at $[[Low]]$ but $[[y]]$ can
only be used when the observer level is $[[High]]$. However, with
\rref{T-AppJoin}, the level for type checking the argument is the
join of the observer level $[[High]]$ and the annotation $[[Low]]$,
so the example holds.

The variant of \rref{T-App} and \rref{T-Abs} that raises the argument
level to be at least high as the observer level can be found in
systems that track either dependency or usage,
such as
QTT~\citep{atkey2018syntax} (after instantiating to the boolean semiring) and
extensions of DCC~\citep{abadi1999core},
including DCC$^{\text{pc}}$~\citep{tse2004translating}, the sealing
calculus~\citep{shikuma2008proving}, and DDC~\citep{choudhury:ddc}.
Notably, the alternative rules still satisfy noninterference.
\ifextended
The intuitive justification is as follows. If
the observer level for an application form is at $[[High]]$, then the
argument will eventually be used in a computation that is allowed to
depend on $[[High]]$ variables. The indistinguishability relation at
level $[[High]]$ will be able to observe the argument regardless of whether the
argument is marked as $[[High]]$ or $[[Low]]$. Therefore, it is safe
to relax the constraint and type check the argument at $[[High]]$
instead. In other words, if noninterference is the only property we
desire from our type system, our more restrictive rules may appear to
be an overkill.
\fi

In earlier iterations of \name{}, we opted for \rref{T-AppJoin} and \rref{T-AbsJoin}.
However, we found that these rules were incompatible
with the use of indistinguishability for conversion,
since the system no longer
admitted the subsumption property.
Consider the following function $f$ that type checks
with both the alternative and existing rules.
\[ [[ empty |- Low \Low x : (Bool High -> type star) . \Low y : ((x High true) Low -> Bool) . \Low z : (x High false) . y Low z : Bool ]]  \]
In the example above, since $[[x High false]]$ and $[[x High true]]$ are
indistinguishable at level $[[Low]]$,
we are allowed to pass $[[z]]$, a term with
type $[[x High false]]$, as an argument to $[[y]]$, a function which
expects something with type $[[x High true]]$.

If subsumption holds, we can type check $f$ at level $[[High]]$.
However, we run into a problem when we try to apply $f$ using
\rref{T-AppJoin}: the first argument for $f$ is type checked
at $[[High \/ Low]]$, where $[[High]]$ is the observer level
of $f$ after applying subsumption and $[[Low]]$ is the annotation from the lambda. That is, the
alternative type system would accept any term $[[a]]$ as an argument
for $f$
as long as $[[empty |- High a : Bool High -> type star]]$.
But since $[[a]]$ is typed at $[[High]]$, it can
produce different types by pattern matching against its
input. By instantiating $[[a]]$ properly, $f$ would allow us
to convert between two arbitrary types!
On the other hand, with our existing rules, we enforce that $[[a]]$ is
well typed at $[[Low]]$ regardless of the observer level, so
$[[a]]$ will always behave
like a constant function even after applying subsumption.

The above contradiction can also be phrased in terms of the
incompatibility between indistinguishability and the
following, which we informally refer to as the upgrade
property.

\begin{proposition}[Upgrade property, does not hold for \name{}]
\label{prop:upgrade}
If $[[W |- psi b : A]]$, then $[[psi \/ W |- psi b : A]]$, where
$[[psi \/ W]]$ is an operation that joins the levels in $[[W]]$
pointwise against the level $[[psi]]$.
\end{proposition}

The proposition states that raising the levels in the context to the
observer level does not make the term $[[b]]$ more difficult to type
check. Suppose we have $[[x : Low Bool High -> type star]]$ in
the typing context $[[W]]$. Then, the derivation of $[[b]]$ could
convert between $[[x High false]]$ and $[[x High true]]$ via
indistinguishability. By applying the upgrade property, however,
we end up with $[[ x : High Bool High -> type star ]]$ and $[[x]]$ can no
longer be used at observer level $[[Low]]$ in that derivation.
As a result, the upgrade property does not hold
in \name{}.

In both QTT and 
DDC, we can find lemmas that capture the same idea as
the upgrade property. In QTT, the upgrade property holds when
its semiring structure is instantiated to $\{0,1\}$, the boolean semiring. In DDC, the
upgrade property is more interesting. 
DDC is parameterized by
some lattice structure $\mathcal{L}$ and additionally adds a two-point
lattice $\{C, \top\}$ on top of $\mathcal{L}$ such that $[[psi]] < C < \top$
for all $[[psi]] \in \mathcal{L}$. Like \name{}, DDC uses
indistinguishability as part of its conversion rule; however, the
observer level for indistinguishability in the conversion rule is
always $C$. In fact, DDC treats the level $\{C, \top\}$ differently
compared to the levels from the lattice $\mathcal{L}$.

The upgrade
property in DDC only holds when the level $[[psi]]$ is less than or
equal to $C$. The disjoint treatment between $\{C, \top\}$ and
$\mathcal{L}$ in DDC is consistent with our finding: since the lower
levels satisfy the upgrade property, indistinguishability can no
longer be used for conversion at those levels. In contrast,
by abandoning \rref{T-AppJoin, T-AbsJoin} and therefore the upgrade property,
\name{} can use indistinguishability for conversion at any level in the lattice.

\subsection{Order Constraints for Indexed Equality Types}
\label{sec:orderconstraints}

We explain the necessity of the ordering constraint
$[[psi0 <= psi]]$ that appears in \rref{GE-CEq}. % Note
% that \rref{T-CEq} additionally has the constraint $[[psi1 <= psi0]]$,
% which has already been explained in Section~\ref{sec:equality}.
\[ % \drule[width=3in]{T-CEq} \quad
  \drule[width=4.0in]{GE-CEq} \]
In \rref{GE-CEq},
to show that the two equality types $[[psi0 a0 ~ b0 : psi1 A0]]$ and
$[[psi0 a1 ~ b1 : psi1 A1]]$ are equal, the endpoints $[[a0]]$, $[[a1]]$
and $[[b0]]$, $[[b1]]$ are compared pointwise at the observer level $[[psi]]$,
where $[[psi]]$ is constrained to be greater than or equal to the
observer level $[[psi0]]$ of the equality type. The premise of
\rref{GE-CEq} does not require any relation between $[[A0]]$ and
$[[A1]]$ since they are type annotations.

It is tempting to
simplify \rref{GE-CEq} into the following.
\[\drule[width=4.0in]{GE-CEqWrong}\]
Compared to \rref{GE-CEq}, this simplified version completely ignores the
observer level $[[psi]]$ % attached to the equality type 
in the premises and instead uses the level $[[psi0]]$ to compare the subterms.
The issue
with this design is that it breaks the congruence property of
indistinguishability (Lemma~\ref{lemma:gdefeqcong}).
Consider the following derivable judgment with \rref{GE-CEqWrong}.
\[    [[x : High |- Low (High x ~ true : High Bool) == (High
    x ~ true : High Bool) ]] \]
Since $[[not (High <= Low)]]$, we can derive the vacuously true guarded
indistinguishability $[[empty |- Low High false == true]]$. If
Lemma~\ref{lemma:gdefeqcong} holds, then we can apply substitution and
derive this indistinguishability judgment.
\[ [[empty |- Low (High false ~ true : High Bool) == (High true ~ true : High
  Bool)]] \]
The equality on the left is bogus, but the equality on the right is
true by reflexivity. We could thus
inhabit the equality type $[[High false ~ true : High Bool]]$
with $[[refl]]$ by composing \rref{T-Refl} and \rref{T-Conv} with
the indistinguishability judgment from above. By eliminating the bogus proof term, we could easily
derive a crashing program.

The problem with \rref{GE-CEqWrong} is that
the two equality types have their own notion of indistinguishability
based on the level $[[psi0]]$, which is now independent of the observer
level $[[psi]]$. The congruence lemma allows us
to substitute in an equality that is sensible for an observer at level $[[psi]]$.
However, if $[[psi]] < [[psi0]]$, then the equality being
substituted in is too coarse for an observer at level $[[psi0]]$. In
the counterexample above, the variable $[[x]]$ is labeled with level
$[[High]]$. From an $[[Low]]$ observer's perspective, $[[false]]$ and
$[[true]]$ are indistinguishable when guarded by the level
$[[High]]$. However, from the $[[High]]$ observer's perspective,
$[[false]]$ and $[[true]]$ are distinct when guarded by the level
$[[High]]$. The mismatch between the notion of equivalence allows us
to break type soundness in a way similar to the $\ottkw{unbox}$
example from Section~\ref{sec:equality}.

To address this problem, \rref{GE-CEq} restricts the observer level of
the equality type to be at most the observer level of the
judgment.
The same constraint $[[psi0 <= psi]]$ found in
\rref{T-CEq} is a direct consequence of the constraint in
\rref{GE-CEq} since we need the order constraint in the typing rule to
derive Lemma~\ref{lemma:typinggdefeq}.

\subsection{Irrelevance, Relatively}
\label{sec:irrelrel}
To talk about whether two
terms are indistinguishable, we must choose an observer level.
Terms that are indistinguishable at $[[Low]]$ might be
distinguishable at $[[High]]$.
So, irrelevance in \name{} is
relative. For example, if $[[empty |- Low a : Bool High -> T High Bool]]$, then these terms are both valid candidates for $[[a]]$.
\[ [[\High x : Bool . box High x]]  \qquad
  [[\High x : Bool . box High false]] \]
The second function behaves like a constant function in
the traditional sense---the input $[[x]]$ does not appear
in the body. The first function is only a
constant function if we observe its results at level $[[Low]]$.

Consider the \lstinlinecol{filter} function on vectors as another example,
where unmarked arguments and the function itself are at level $[[Low]]$.
Because we do not know the final length of the filtered vector,
we package the vector and its length together in a dependent pair
where the first component is at level $[[High]]$,
which we write as \lstinline{(m :$^[[High]]$ Nat) $\times$ Vec A m}.

\begin{lstlisting}
    filter : (A :$^[[High]]$ Type) -> (n :$^[[High]]$ Nat) -> (A -> Bool) -> Vector A n ->
             (m :$^[[High]]$ Nat) $\times$ Vector A m
    filter A n f v = case v of
      Nil => (0, Nil)
      Cons n' y ys => let zs = filter A n' f ys in
                     if f y then (Succ (fst zs), Cons (fst zs) y (snd zs)) else zs
\end{lstlisting}

From the perspective of a run-time observer,
the arguments \lstinline{A} and \lstinline{n} are both irrelevant.
However, for an observer at level $[[High]]$, these arguments are both relevant.
In the body of \lstinlinecol{filter},
the variable \lstinline{zs} is a pair containing a
vector of unknown length. The length index \lstinline{fst zs} in the
pair is irrelevant at level $[[Low]]$; after erasure at
level $[[Low]]$,
\lstinlinecol{filter} should have the run-time performance of the
\lstinlinecol{filter} function for lists.
On the other hand, the length index in the pair is used relevantly
in the type of the second component, \lstinline{Vector A (fst zs)}.
Because the notion of relevance is relative, we
can still make sense of the \lstinlinecol{filter} function even though
the length index is used both relevantly and irrelevantly.

\ifextended
\subsection{Indistinguishability and Type Conversion}
\scw{Because this alternative does not actually change the language, if we are short 
on space, we can skip this section. It's mostly an observation that we want definitional 
equality to be transitive.}
\Rref{T-Conv} can alternatively be
formulated in terms of indistinguishability directly.
\[\drule[width=3.0in]{T-ConvAlt}\]
Given $[[P |- a == b]]$, it is not always possible to find some level
$[[psi]]$ such that $[[P |- psi a == b]]$. It turns out
this alternative design is not more restrictive than the
\rref{T-Conv} since a single derivation of \rref{T-Conv} can be
decomposed into multiple applications of \rref{T-ConvAlt}. 

However, in Section~\ref{sec:preservation}, the statement of the
inversion lemma relies on definitional equality. For example, given 
$[[W |- psi \psi0 x : A . a : A0]]$, \rref{T-ConvAlt} could
be applied multiple times to convert from the function type to
$[[A0]]$. Each application of \rref{T-ConvAlt} may apply
indistinguishability at a different level and therefore we still need
to rely on definitional equality to relate the function type and
$[[A0]]$. We choose to use definitional equality in \rref{T-Conv}
to be more consistent with the formulation of PTS, which uses
$\beta$-equivalence for both the conversion rule and the statement of the
inversion lemmas.

\scw{We could move this example to the explanation of $[[ P |- a == b ]]$. I think it 
helps explain the expressiveness of the relationship.}
\yl{I find the example  too pathological to help explain
  expressiveness. That's why I moved it here from Section 3.}

Finally, we note that given two definitionally equal terms $[[a]]$ and $[[b]]$,
it is not always possible to find some level $[[psi]]$ such that
$[[a]]$ and $[[b]]$ are indistinguishable at $[[psi]]$.
\begin{align*}
&[[y : Low |- High (unbox High (box High y)) High false   == y High true ]]\\
&[[y : Low |- Low y High false == y High true ]]
\end{align*}
The definitional equality  $[[y : Low |- (unbox High (box High y)) High false == y High true]] $
holds by chaining the two indistinguishability judgments above.
However, these two terms cannot be equated through
indistinguishability at either $[[Low]]$ or $[[High]]$.

\yl{Removed some text here because they were meant to be intuitive
explanations of the conversion rule but I don't find my explanation
intuitive after all..}

\yl{deleted the text here because there's a section that already
  discusses relative irrelevance }
\fi

\subsection{Normalization Proof}
\label{sec:normalization}

In PTS, particular instantiations of the sorts, axioms, and rules lead
to systems that are known to satisfy the strong normalization
property.
We have not proven normalization for any instance of \name{};
however, since it is based on PTS, we believe that there is a close
correspondence between the normalization behavior of PTS and of \name{}.
If a specific instantiation of the sorts results in a
normalizing PTS, we would expect the same instantiation to result in a
normalizing \name{}. In particular, we expect all instances from the
lambda cube~\citep{barendregt1991introduction} to satisfy the strong normalization property when
reformulated as their corresponding \name{} instance.

Having a normalization proof for such instantiations of sorts is
desirable to allow for relevance tracking with proof erasure.
In \rref{T-Transp},
the level of the equality proof must be the same as the level of the term being cast,
so we run into the awkward situation where an
irrelevant proof cannot be used to cast a relevant term. In a system
where equality types can be inhabited by a diverging computation, the
proof being eliminated must be treated as relevant since the erasure
of a bogus proof can cause the resulting program to crash.
However, if we know that a particular instance of \name{} is normalizing,
then we are able to justify the elimination of an irrelevant equality proof in a
relevant context through the canonicity of closed equality proofs.
Because of its importance, we will investigate the normalization behavior of \name{} as
part of our future work.

\ifextended
Furthermore, in a normalizing system, we are allowed to
define an extensional version of the indexed equality
type. Unlike intensional equality, extensional equality does not
induce any run-time overhead since an equality proof can be directly
reflected into definitional equality without having an elimination
form in the term syntax. In instantiations of \name{} that do not
normalize (e.g. with the axiom $\{(\ast, \ast)\}$),
it would be possible to convert between arbitrary types by
reflecting bogus proofs, and we would no longer be able to show type
soundness with an extensional equality type.
\fi

\section{Implementation}
\label{sec:implementation}

We have implemented a type checker for an instantiation
of \name{} with the naturals and its usual total order as our lattice of labels,
and a single type universe \lstinline{Type} with type-in-type\punctstack{.}%
\footnote{The code in this section has been stylized to match existing examples;
the actual concrete syntax is introduced in \dotpi{impl}{README.pi}.}
The biggest difference between \name{} and the implementation is that
the implementation contains data types,
which subsumes the unit type and dependent pair types.
Below is a data definition for dependent pairs
where the first component is fixed to level 1;
we use it to discuss some design decisions for data definitions.

\begin{lstlisting}
    data Pair (A :$^0$ Type) (B :$^0$ A$^1$ ->  Type) :$^0$ Type where
      MkPair :$^k$ (x :$^1$ A) -> (B x)$^k$ -> Pair A B
\end{lstlisting}

The level of type parameters \lstinline{A} and \lstinline{B}
and of the \lstinline{Pair A B} type are fixed to level 0,
in contrast to the primitive dependent pair type of \name{},
which are implicitly level-polymorphic in the sense that they can be
type checked at any level
as long as the parameters can be type checked at the same level.
Instead of fixing to level 0, we could assign \lstinline{A} and \lstinline{B} different levels,
as long as the level of \lstinline{Pair A B} is greater than or equal to
the larger of these levels. This condition ensures that the parameters are always relevant:
we should not to equate \lstinline{Pair A$_1$ B$_1$} and \lstinline{Pair A$_2$ B$_2$}
unless \lstinline{A$_1$}, \lstinline{A$_2$} and \lstinline{B$_1$}, \lstinline{B$_2$} are themselves equal.
\scw{Is there a sound extension that allows irrelevant parameters to datatypes?}
\jc{Seems unlikely. The parameter would probably have to never be used later,
but if that were the case, why would the parameter exist at all?
If there's a different scenario where it's irrelevant but also used in a valid way,
we could speculate in the extensions section.}
To retain the subsumption property, data types can be used at any level
greater than or equal to the level they are declared at,
similar to how variables are checked by \rref{T-Var}.

On the other hand, constructors are explicitly level-polymorphic:
the second component of the pair is type checked at the same level variable $k$ as the overall pair,
so that when given a pair declared at a given level,
we can use the second component at the same level.
This minimal amount of polymorphism appears to be sufficient for practical examples.
As with \rref{T-SPair}, there is no relation between the level of pair type and $k$.

Pairs are eliminated by case expressions rather than projections,
but projection functions can still be defined.

\begin{lstlisting}
    proj1 :$^1$ (A :$^0$ Type) -> (B :$^0$ A$^1$  -> Type) -> (Pair A B)$^0$ ->  A
    proj1 = \A B p. case p of
      MkPair x y => x

    proj2 :$^0$ (A :$^0$ Type) -> (B :$^0$ A$^1$  -> Type) -> (p :$^0$ Pair A B) -> B (proj1 A B p)
    proj2 = \A B p. case p of
      MkPair x y => y
\end{lstlisting}

The branches of a case expression must be checked at a level greater or equal to that of the scrutinee.
Intuitively, this prevents extracting relevant information out of something irrelevant.
As a minor detail, in the type of the branch for a given constructor,
the scrutinee is replaced by the constructor pattern,
so the branch of \lstinline{proj2} requires a term of type \lstinline{B (proj1 A B (MkPair x y))},
reducing to \lstinline{B x}, which is exactly the type of \lstinline{y}.

If we have pairs at level 0, then as expected,
two pairs are indistinguishable at level 0 merely if their second components are,
since their first components (at level 1) are automatically indistinguishable at this level.
With internalized indistinguishability, this can be proven as a lemma.

\begin{lstlisting}
    pairEq :$^0$ (A :$^0$ Type) -> (B :$^0$ A$^1$ -> Type) -> (p1 :$^0$ Pair A B) -> (p2 :$^0$ Pair A B) ->
             (proj2 A B p1 ~$^0$ proj2 A B p2) -> p1 ~$^0$ p2
    pairEq = \A B p1 p2 q. case p1 of
      MkPair x1 y1 => case p2 of
        MkPair x2 y2 => transp q refl
\end{lstlisting}

In \name{}, definitional equality is the transitive closure of indistinguishability
at existentially-quantified levels; in the implementation, when converting between two types,
we simply use indistinguishability at an arbitrary generated level metavariable
at which both types are well typed.

The conversion checker checks guarded indistinguishability $[[P |- psi psi0 a == b]]$
by first asking the constraints whether $[[psi < psi0]]$ is implied.
If so, this corresponds to the \rref{CGE-NLeq} case;
otherwise, it adds $[[psi0 <= psi]]$ as a new constraint
and continues on as in the \rref{CGE-Leq} case.
Conversion checking is therefore incomplete with respect to the rules,
since it may turn out that $[[psi < psi0]]$ was the appropriate constraint to add.
Adding backtracking might make conversion complete,
but it seems that with explicit level annotations,
enough of the guesswork is eliminated.

To accommodate level metavariables, type checking always takes a level as an input,
collecting constraints for generated level metavariables,
and solves the level constraints with minimal level instantiations at the end,
or fails with a type error.

\scw{What is the largest example that you have checked with the implementation?}

\section{Future work}
\label{sec:extensions}

\paragraph{Level polymorphism}
There are multiple forms of level polymorphism that could be added to \name{}.
The simplest is prenex polymorphism which,
in combination with global or let-bound definitions, can reduce code duplication.
By quantifying a definition over the levels it uses,
it can be instantiated at different levels as needed.
We might however wish to also enforce constraints among the levels quantified,
so that we have guarantees about their relative relevancies.
This can be done by bounded quantification,
which can only be instantiated by levels bounded above by the one or more other levels given.
Adding such forms of polymorphism likely would not violate any of our metatheoretical properties,
since each instantiation of a polymorphic definition with particular levels could be inlined
as an ordinary \name{} term with those levels substituted in.

\ifextended
On the other hand, it isn't evident how higher-rank level polymorphism might work,
where functions can abstract over level-polymorphic terms.
In particular, if the relevance of a level-polymorphic term depends
the levels it's instantiated with,
it's unclear what the overall relevance of a function abstracting over it should be.
Additionally, the usefulness of higher-rank level polymorphism needs further investigation.

Finally, we might go as far as first-class levels that can depend on terms.
Such levels could have applications in secure information flow,
where the observer's security access level is only known dynamically at run time.
\fi

\paragraph{Inductive types}

While our prototype implementation contains data types with levels,
adding inductive types to \name{} with proper type indices, strict positivity checking, and eliminators,
and updating the Coq development to ensure none of the metatheoretical properties are broken,
remain future work.
Aside from increasing expressivity,
another goal of adding inductive types is to subsume the existing unit,
dependent pair, and propositional equality primitive types.
Although we won't know what rules for inductive types are correct until the proofs are done,
we can use the implementation's type checker and the rules for the existing types
to guide the design of expressive yet well-behaved inductive types.
\ifextended
For instance, consider the following proposal for propositional equality as an inductive type.

\begin{lstlisting}
    data Eq (A :$^{[[psi2]]}$ Type) (x :$^{[[psi1]]}$ A) :$^{[[psi]]}$ A$^{[[psi1]]}$ -> Type where
      Refl :$^{[[psi3]]}$ Eq A x x @ $[[psi0]]$
\end{lstlisting}

$[[psi]]$ is the level of the fully-applied inductive type,
$[[psi0]]$ the observer level of the equality,
$[[psi1]]$ the level of the equated terms,
$[[psi2]]$ the level of the type of the equated terms,
and $[[psi3]]$ the level of the reflexive proof.
\Rref{T-CEq} enforces $[[psi1]] \leq [[psi0]] \leq [[psi]]$.
In general, we want to make sure that inductive types applied to different parameters and indices
are always distinguishable, so we would additionally enforce $[[psi2 <= psi]]$ like in the implementation.
Next, we might type a case expression for propositional equality as follows.

\[ \drule[width=6.5in]{T-CaseEq} \]

\Rref{T-Transp} enforces $[[psi3 = psi4]]$.
As discussed in Section~\ref{sec:normalization},
in nonnormalizing instances of \name{},
allowing $[[psi4 < psi3]]$ would be invalid,
since this would incorrectly leak the normalization behavior of the higher-level $c$ to a lower level.
However, $[[psi3 <= psi4]]$ may be valid,
and is currently in the implementation,
since it allows destructing an inductive without having to cast its level up first.
\fi
Lastly, inductive definitions should be prenex-polymorphic in all levels mentioned
to be as flexible as the primitives, which are not fixed in their levels.

\paragraph{Staged programming\protect\footnote{\dotpi{impl}{pi/Staged.pi}}}

A potential application of \name{} is in \textit{staged programming}~\citep{sheard1995staged},
where different pieces of code are marked with different levels or \emph{stages},
which can be thought of as different stages of metaprogramming.
Our reduction rules could be extended by indexing with a stage
and modified to proceed stage by stage,
corresponding to expanding metaprograms at each stage in order.

Although the dynamic behavior does not yet correspond to staged evaluation,
we can use the existing static types for staged programs.
In particular, the box construct represents a way to manipulate code at different stages;
for example, a value of type \lstinline{$T^1$ ((Nat$^1$) -> Nat)} is a function over naturals available at stage 1
but is manipulable as an opaque piece of object code at stage 0.

Consider the following implementation of exponentiation of naturals at stage 0
in terms of multiplication at stage 1.
In other words, exponentiation resembles a macro that expands to applications of multiplication.

\begin{lstlisting}
    mul :$^1$ Nat$^1$ -> Nat$^1$ ->  Nat
    exp :$^0$ Nat$^1$ ->  Nat$^0$ ->  $T^1$ Nat
    exp b e = case e of
      Zero => $\ottkw{box}^1$ 1
      Succ e' => $\ottkw{box}^1$ (mul b ($\ottkw{unbox}^1$ (exp b e')))
\end{lstlisting}

With proper staged evaluation,
when we reach stage 1, expressions at stage 0 will have been evaluated,
and applications of exponentiation to a concrete exponent will have been expanded
into iterated multiplication.
For instance, \lstinline{($\ottkw{unbox}^1$ (exp b 3))}
evaluates to \lstinline{(mul b (mul b (mul b 1)))} at stage 1.

Analogously to previous applications of dependency analysis,
pieces of code from later stages are indistinguishable from code from earlier stages.
Intuitively, this means that metaprograms can't tell what programs are doing,
since programs can't run before the metaprograms expand.

Aside from separating programs into different stages of evaluation,
levels could also separate proofs into different type theories,
as is done in two-level type theory~\citep{annenkov2023two},
which also has applications in staged compilation~\citep{kovacs2022staged}.

\begin{comment}
This is provably true within the language as well, thanks to the propositional equality type,
which allows us to reason about code in later stages.

\begin{lstlisting}
    cubed :$^0$ (b : Nat$^1$) -> cube b ~$^1$ mul b (mul b (mul b 1))
    cubed b = $[[refl]]$
\end{lstlisting}

Inductive types can be staged as,
for instance encoding the definition of length-indexed vectors at stage 0
in terms of pairs at stage 1.
Here, we use the nondependent pair type $A \times B$,
where $A$ and $B$ are both also at stage 1,
and arguments without stage annotations are at stage 0 by default.

\yl{What are the benefits of staging a Vector type? Is the idea to have
  staged computation for type-level reduction so type-checking becomes
  faster?}
\jc{I'm not sure. the 2LTT paper included it as an example so it seemed important.
They were demonstrating that it was \emph{possible} to stage types as well as terms.
(They did have a typo in their example, though, missing an unboxing of Vector
in the types of nil and cons.)}
\begin{lstlisting}
    Vector :$^0$ Type$^1$ -> Nat -> $T^1$ Type
    Vector A n = case n of
      Zero => $\ottkw{box}^1$ $[[Unit]]$
      Succ n' => $\ottkw{box}^1$ (A $\times$ $\ottkw{unbox}^1$ (Vector A n'))

    nil :$^0$ (A :$^1$ Type) -> $\ottkw{unbox}^1$ (Vector A 0)
    nil A = $[[unit]]$

    cons :$^0$ (A :$^1$ Type) -> (n : Nat) -> A ->
            $\ottkw{unbox}^1$ (Vector A n) -> $\ottkw{unbox}^1$ (Vector A (Succ n))
    cons A n a v = (a, v)

    fold :$^0$ (A :$^1$ Type) -> (P : (n : Nat) -> $\ottkw{unbox}^1$ (Vector A n) -> Type) -> P 0 (nil A) ->
           ((n : Nat) -> (a : A) -> (v : $\ottkw{unbox}^1$ (Vector A n)) ->
            P n v -> P (Succ n) (cons A n a v)) ->
           (n : Nat) -> (v : $\ottkw{unbox}^1$ (Vector A n)) -> P n v
    fold A P pnil pcons n v = case n of
      Zero => pnil
      Succ n' => pcons n' (proj1 v) (proj2 v) (fold A P pnil pcons n' (proj2 v))
\end{lstlisting}

Code at stage 0 interacts with vectors through this interface,
while applications to vectors of concrete lengths get compiled away to functions on pairs at stage 1.
For example, \lstinline{$\ottkw{unbox}^1$ Vector A 2} becomes \lstinline{A $\times$ A $\times$ $[[Unit]]$},
\lstinline{cons A 1 a$_1$ (cons A 0 a$_0$ (nil A))} becomes \lstinline{(a$_1$, a$_0$, $[[unit]]$)},
and \lstinline{fold A P pnil pcons 2 v} becomes \lstinline{pcons 1 (fst v) (snd v) (pcons 0 (fst (snd v)) (snd (snd v)) pnil)}.
These equalities, too, are provable propositionally.

\jc{I'm reluctant to mention 2LTT (HoTT + MLTT) at this point.
All of the references I can find talk about how HoTT is the ``inner'' type theory
and MLTT is the ``outer'' type theory. The Annenkov paper also says:
\begin{quote}
  Thus, we can convert inner types into outer one, but not always vice versa.
\end{quote}
But this is the \emph{opposite} of how our system works:
level 0 is the metalevel, level 1 is the object level,
and by subsumption, anything typeable at level 0 can also be typeable at level 1,
meaning it should be possible to go from MLTT (the meta) to HoTT (the object),
but that's exactly the opposite situation here.
(It's certainly not valid to lift axiom K into HoTT!)
None of the papers use a secret third type theory that's the meet of MLTT and HoTT either,
it's always an inner HoTT and an outer MLTT.
I don't want to speculate on something unconventional in the examples;
maybe it's better suited for the discussion,
but even then I don't know if it's worth saying many things incorrectly
over saying fewer things correctly.
Maybe reviewers would take more kindly to suggesting a missing connection
than to correct a blantant misunderstanding.
And if one of them \emph{does} point out the connection to 2LTT,
then we'd know with greater certainty that the speculation was on track.}
\end{comment}

\section{Related work}
\label{sec:related}

\subsection{Systems with Dependency Tracking}
The Dependency Core Calculus (DCC)~\citep{abadi1999core} introduces the
graded model type $[[T psi A]]$ for dependency tracking. Similar to the box
type from \name{}, given a term $[[a]]$ of type $[[T psi A]]$, the
content stored inside $[[a]]$ can only be used in a context that
can observe $[[psi]]$. DCC enforces
information flow by restricting the type of the \lstinline{bind}
operator that inspects the content of $[[a]]$. For example, supposing
$[[a]]$ has type $[[T psi0 Bool]]$, the return type of the bind
must then be protected at level $[[psi0]]$. Intuitively, a type is
protected at $[[psi0]]$ when every branch node of the type's syntax
tree is guarded by some $T^{[[psi]]}$ constructor where $[[psi0 <=
psi]]$, excluding domain types.

Variations of DCC, such as DCC$^{\text{pc}}$ from
\citet{tse2004translating} and the sealing
calculus from \citet{shikuma2008proving}, introduce a program counter
or observer level to the typing judgment. Accessing the data
stored in a box of type $[[T psi0 A]]$ is done by comparing the observer
level $[[psi]]$ against the level of the box $[[psi0]]$, similar to
the admissible \rref{T-Unbox} in \name{}. 

The Dependent Dependency Calculus (DDC)~\citep{choudhury:ddc} extends
PTS~\citep{barendregt:lambda-calculi-with-types} with dependency tracking.
Compared to previous systems,
DDC's design was influenced by systems with coeffects, and the
variables in the typing context are labeled with levels.
The lattice $\mathcal{L}$ in DDC is augmented with two elements $\{C, T\}$ such that
$[[psi]] \leq C \leq T$ for every $[[psi]] \in \mathcal{L}$.
The dependency tracking mechanism for levels
below $C$ is similar to previous systems,
whereas the levels $C$ and $T$
are handled differently through the 
resurrection mechanism~\citep{pfenning:irrelevance} in order to support
compile-time irrelevance. The type conversion rule of DDC is based on
indistinguishability, though the observer level of the
indistinguishability judgment must be fixed to $C$, due to the
limitation discussed in Section~\ref{sec:upgrade}.
The design of DDC directly inspired the use of indistinguishability in
\name{}'s conversion rule. However, this work handles dependency tracking and
compile-time irrelevance through a uniform mechanism, generalizes
the conversion rule to arbitrary observer levels, and includes propositional equality.

\subsection{Compile-Time Irrelevance}
The ability to convert between types based on indistinguishability can
be viewed as a generalization of a form of compile-time
irrelevance. This feature plays an important role in
reasoning about data structures containing embedded proofs: the type checker
should treat different witnesses of the same proposition as equal and
only compare relevant components.

In Coq and Agda, the respective sorts \lstinline{SProp} and
\lstinline{Prop}~\citep{gilbert2019definitional} classify types whose
inhabitants are treated as definitionally equal. The system then ensures that
terms whose types are in this sort can never be pattern matched to produce
terms of other sorts, preventing the flow of information from irrelevant types
to computationally relevant types. The \lstinline{Squash} constructor
allows one to embed data types of the relevant sort into
\lstinline{SProp} and is analogous to the graded modal type $[[T psi
A]]$ from \name{}.
In \citet{gilbert2019definitional}, it is
impossible to construct the choice function with the
signature \lstinline{Squash (Squash A -> A)}. Instead, the choice
function most be postulated as an axiom.
In \name{}, the choice function with the analogous type $[[T High (Bool
High -> Bool)]]$ can be defined as the term $[[ box High (\High x : Bool
. x)]]$.

% DCOI, on the other hand, does not
% associate sorts with relevance and is more closely related to
% \citet{pfenning:irrelevance} and \citet{Abel12}, which implement proof
% irrelevance with irrelevant function types.

Resurrection, introduced by \citet{pfenning:irrelevance}, is used by DDC,
EPTS~\citep{mishra2008erasure}, and Agda~\citep{Abel12}. Similar to
\name{} and unlike \citet{gilbert2019definitional}, proof irrelevance
in \citet{pfenning:irrelevance}
is not determined by sorts.
As in \name{}, the type system
can label variables as either relevant or irrelevant. However, a type system
based on resurrection does not require an observer level in its typing
judgment. Instead, when checking irrelevant parts of the term, the type system
uses the resurrection operation on the typing context to make irrelevant
variables accessible.

To see how resurrection behaves differently from the
dependency tracking mechanism of \name{}, consider the terms
$ [[ \High x : Bool . box High x]]$ and $[[ box High (\High x : Bool
. x)]]$, the latter of which is the choice function we have seen earlier.
In \name{}, both terms type check at level $[[Low]]$. However, with
resurrection, only the first term type checks. 
% To understand why, let
% us first see how the first type checks through resurrection.
In the first term,
by the
abstraction rule, it suffices to show that $[[box High x]]$ is well typed given the typing context $[[x : High
Bool]]$. Before checking the subterm wrapped inside an $[[High]]$ box,
the typing context is resurrected and the level of the variable $[[x]]$ is lowered
to $[[L]]$. As a result, the $[[x]]$ can legally appear inside the
box.
%
In the second term, the $[[High]]$ box wraps around
the entire lambda expression. Before checking $[[(\High x : Bool
. x)]]$, the empty context is resurrected. In this case, resurrection
has no effect on the variable $[[x]]$, which has yet to be introduced
to the typing context, so
its usage in the
body of the lambda is rejected.

Resurrection interacts poorly with strong dependent pairs because
projections, unlike pattern matching, do not introduce
variables that can be resurrected later on.
Consider a term \lstinline{xs} of type
\lstinline{(m :$^[[High]]$ Nat) $\times$ Vector A m}. %When checking
% terms, the resurrection mechanism correctly rejects
% \lstinline{proj1 m} since the first component of \lstinline{as} is
% irrelevant and should be erased.
We cannot check the second projection \lstinline{proj2 xs}
because its type, \lstinline{Vector A (proj1 xs)}, contains the
ill-typed term \lstinline{proj1 xs}. If we interpret the level
$[[High]]$ as describing irrelevance at runtime, 
rejecting \lstinline{proj1 xs} in the term but accepting
\lstinline{proj1 xs} in the type is the desired behavior. 
However, without an observer level on the typing judgment,
there is no way to distinguish
between the illegal use of \lstinline{proj1 xs} in the term and
the legal use of \lstinline{proj1 xs} as an argument to
\lstinline{Vector}. Both must be
rejected.

As pointed out by \citet{choudhury:ddc}, a direct consequence of the above
limitation is that the \textsf{filter} function from
Section~\ref{sec:irrelrel} must pattern match the returned existential from
the recursive call. \citet{eisenberg:existentials} observe that
\textsf{filter} defined using pattern matching is needlessly strict and loses
the streaming behavior of an ordinary list filter function in a
call-by-name language.

To overcome the limitations of the resurrection mechanism, DDC restricts labels
for run-time irrelevance from being used as the observer in definitional
equality. In Agda, there is a disjoint mechanism for run-time irrelevance and
function arguments are considered independently for run-time and compile-time
irrelevance. \Name{} provides a uniform mechanism that works the same for
reasoning about run-time erasure and compile-time equality.

% https://agda.readthedocs.io/en/v2.6.3/language/runtime-irrelevance.html
% https://agda.readthedocs.io/en/v2.6.3/language/irrelevance.html
% https://agda.github.io/agda/Agda-TypeChecking-Irrelevance.html

\subsection{Quantitative Type Systems}
Unlike systems that track only what depends on what, quantitative type systems
track usage more precisely. In these systems, inputs are
annotated with abstract usage information drawn from an arbitrary
semiring. This 
information specifies how many times a variable may be
used in the computation. These systems are related to dependency tracking because
computation may not depend on a variable marked with $0$ usage.

The run-time irrelevance mechanism in Agda~\citep{abel:icfp2020} and Idris~\citep{brady:idris2}
is based on Quantitative Type Theory (QTT)~\citep{mcbride2016got,atkey2018syntax}.
The typing judgment in QTT takes the form $\Gamma \vdash M:^\sigma S$
where $\sigma \in \{0,1\}$.
The $\sigma$ at the colon indicates the relevance of the term being
constructed, similar to the observer level in \name{}.
When $\sigma = 1$, the term $M$ is computationally relevant.
When $\sigma = 0$, the term $M$ is irrelevant.
A variable with $0$ usage can only be used when
constructing irrelevant terms or in types.

While QTT judgments resemble those of \name{}, QTT only tracks usage at the
term level: in types, all tracking is effectively disabled.  As a result,
QTT's definitional equality cannot rely on usage information.  Furthermore, as
mentioned in Section~\ref{sec:upgrade}, when instantiated with the boolean
semiring, QTT satisfies the upgrade property (Proposition~\ref{prop:upgrade})
using the $1 \leq 0$ ordering. As a result, QTT cannot replace
$\beta$-equivalence with a usage-aware equality without breaking type
soundness.

% Note that when QTT's semiring is instantiated with the boolean semiring, 
% then 
% \yl{My conjecture is that QTT with a boolean semiring is isomorphic
%   to DDC Top with a two-point lattice.}
% \yl{If you find the following discussion too risky to include, we can
% remove it. We would need to also remember to remove the atkey citation
% in Section~\ref{sec:upgrade}}
%Interestingly, when the semiring $R$ is instantiated to the boolean
%semiring $\{0,1\}$, the multiplication and addition operators
%correspond to join and meet for the two-point lattice $\{0,1\}$ with
%the ordering $1 \leq 0$. 
%
% the following property, which holds in QTT,
% can be interpreted as the upgrade property
% (Proposition~\ref{prop:upgrade}) using the $1 \leq 0$ ordering.
% \begin{proposition}[QTT: Tm-Zero]
%   \label{prop:tmzero}
%  If $\Gamma \vdash M :^\sigma S$, then $0\Gamma\vdash M :^0 S$ where $0
% \Gamma$ is defined as $\Gamma$ with all usages reset to $0$.
% \end{proposition}
% Based on the reasoning in Section~\ref{sec:upgrade},
% Proposition~\ref{prop:tmzero} implies that QTT  % if we were to extend QTT with usage tracking for types.

The more recent systems, GraD~\citep{weirich:graded-haskell} and
GRTT~\citep{moon2021graded}, track usage in both terms and types. However,
both systems require a fixed point of reference; the typing judgment is not
annotated by a level. (One can view the typing judgment as implicitly indexed
by usage $1$, meaning that the term being constructed is always
computationally relevant or ``public''.)

In GraD, the typing judgment takes the form $\Delta; \Gamma \vdash a : A$
where $\Gamma$ is the typing context with usage information and $\Delta$ is
the same as $\Gamma$ but without usage information. GraD's regularity property
(shown below) is reminiscent of the regularity property
(Lemma~\ref{lemma:regularity}) of \name{}.
\begin{lemma}[GraD: Regularity]
If $\Delta; \Gamma \vdash a : A$, then there exists some $\Gamma'$
such that $\Delta;\Gamma' \vdash A : type$.
\end{lemma}
% Note that the judgments from the premise and the conclusion share the
% same $\Delta$: ignoring the usage information, the context used for checking the
% type $[[A]]$ is identical to the context for checking the term $[[a]]$. 
As in \name{}, there is independence between term-level and type-level usage
information. However, in GraD, the independence appears in the context
($\Gamma'$ vs. $\Gamma$), whereas in \name{}, the independence appears in the
observer level of the judgment.

\ifextended
The ability to track usage at the type level suggests that it might be
possible to implement a usage-aware conversion rule in GraD. In fact,
GraD parameterizes over the equivalence relation $A \equiv B$ used for type
conversion and imposes certain axioms on the binary relation to ensure
type soundness.
While an explicit instantiation is not given besides
$\beta$-equivalence, \citet{weirich:graded-haskell} hints that it is
possible to instantiate the relation to take advantage of usage
information and implement a form of compile-time irrelevance, though
such a claim seems unlikely since the equivalence relation is not
parameterized over the context and thus has no way to access the usage information.
\scw{I'm not sure that I would make this claim today. I don't believe GraD can 
  support compile-time irrelevance any more.}
\yl{I put the entire paragraph to extended. When you say ``I don't believe GraD can 
  support compile-time irrelevance any more.'' is it because the relation is binary and is not indexed by the context?}
\fi

GRTT takes a more fine-grained approach to tracking type-level resource
usage. Each variable in GRTT is associated with two usage levels, one for
term-level usage and one for type-level usage. A variable may be restricted to
the $0$ usage at the term level, but still have a nonzero usage in the type,
and vice versa. In \name{}, we allow a function's argument to be used
relevantly in its type but irrelevantly in its body. However, there is no
way to enforce the invariant that a function's argument is never used in
its type since well-formedness checks in \name{} can be done at an
arbitrarily chosen level.
% \scw{I don't understand what this next sentence is trying to
%   say, so I can't update it. However, it should replace ``parametricity
%   properties'' with something more specific---maybe about reasoning about the
%   use of large eliminations.}
% \yl{I'm fine with deleting the comment about parametricity. }
Despite GRTT's fine-grained usage tracking, its conversion rule is 
$\beta$-equivalence. % In \name{}, given $[[W |- Low \High x : A . a : Pi x : High A . B]]$, we know that $[[x]]$ is used
% in irrelevantly in $[[a]]$ at observer level $[[Low]]$, but the type
% $[[Pi x : High A . B]]$ can be independently checked at level
% $[[High]]$. As a result, the dependency tracking mechanism in \name{}
% cannot recover parametricity like GRTT does.

\section{Conclusion}
In this work, we present \name{}, an extension of Pure Type Systems
with dependency tracking that uses indistinguishability for
type conversion and internalizes indistinguishability as an indexed
equality type. These new features allow reasoning about
information flow within the type system,
extending the expressiveness of the language.
Our design also handles run-time and compile-time irrelevance
in a uniform way, using the same mechanism for
generating more efficient code and reducing the type-checking effort for
programs that rely heavily on type-level computations. 

%These new features allow programmers to reason about 
%information flow within the type system, extending the expressiveness
%of the language. This design also handles run-time and compile-time
%irrelevance in a uniform way, allowing us to use the same mechanism
%to generate more efficient code and reduce the compile
%time of programs that rely heavily on type-level computations.

%%% SCW - no: this future work is still to specific and tied to us. 
%%% we have a section already saying what we will do. Future speculation
%%% should only be about what can be built on top of our ideas.
%%% (But we are out of space anyways, so just cutting.)
% As part of our future work, we will explore ways to reduce the
% level annotations to better integrate our relevance tracking mechanism
% into existing proof assistants and dependently typed languages.


\begin{acks} 
The authors would like to thank the anonymous reviewers for their comments and suggestions.
This work was supported by the \grantsponsor{}{National Science Foundation}{} under Grant Nos.~\grantnum{}{2006535} and
\grantnum{}{2327738}.
\end{acks}

\section*{Data Availability Statement}
The Coq proofs and the Haskell prototype implementation are
available on the ACM Digital Library~\citep{dcoiartifact}.


\bibliographystyle{ACM-Reference-Format}
\bibliography{weirich, refs, students}

% \appendix
\scw{TODO (later): revise proofs to only require semilattice, not lattice}
\jws{TODO (later): do you all prefer Section rather than $\S$?}
\jc{TODO (never): we should really be using `\\cref' for this :p}
\jws{TODO (later): I think the names of the defeq rules are still based on the old conventions (e.g. GE, CGE)}

% \renewcommand\ranchor[1]{{\ranchorstyle{#1}}}

% \section{Rules}

% \subsection{Definitional Equality}
% \label{sec:defeqrules}

% \drules[E]{$ [[P |- a == b]] $}{Definitional equality}{One, Trans}
% \drules[GE]{$ [[P |- psi a == b]] $}{Indistinguishability}{Var, Type, Sym, Trans, App, Abs, AppAbs, Pi, TyUnit, TmUnit,
%   SSigma, ProjOneBeta, ProjTwoBeta, SPair, ProjOneCong,
%   ProjTwoCong,CEq,TranspRefl,Transp,Refl, Bool, True, False, If,
%   IfTrue, IfFalse}
% \drules[CGE]{$ [[P |- psi psi0 a == b]] $}{Guarded indistinguishability}{Leq,NLeq}

% \subsection{Reduction}
% \label{sec:red}
% \drules[R]{$[[a ~> b]]$}{Reduction}{App, AppAbs, Transp, TranspRefl,
%   ProjOneCong, ProjTwoCong,  ProjOneBeta, ProjTwoBeta , If, IfTrue, IfFalse}

% \subsection{Parallel Reductions}
% \label{sec:parrules}
% \drules[EP]{$[[ P |- a => b]]$}{Parallel reduction}{One}
% \drules[Par]{$[[P |- psi a => b]]$}{Indexed parallel reduction}{Var,
%   Type, Pi, AppBeta, App, Abs, Refl, Transp, CEq, TranspRefl,
%   ProjOneCong, ProjTwoCong, ProjOneBeta, ProjTwoBeta, SPair, SSigma, Bool, True, False, If,
%   IfTrue, IfFalse}
% \drules[CPar]{$[[P |- psi psi0 a => b]]$}{Guarded parallel
%   reduction}{Leq, Nleq}

\end{document}
\endinput
%%
